% concatenated from main.tex
\documentclass{lmcs} %%% last changed 2014-08-20

%% optional lists of keywords
\keywords{Solvability Complexity Index (SCI); computable analysis; Weihrauch reducibility; iterated limits; Borel hierarchy}

\usepackage{amsfonts} 
\usepackage{amssymb}
\usepackage{amsthm}
\usepackage{amsmath} 
\usepackage{mathrsfs}
\usepackage{mathtools}
\usepackage{dsfont}
\usepackage{esint}
\usepackage{ulem}
\usepackage{marginnote}
\usepackage{array}
\usepackage{algorithm}
\usepackage{algpseudocode}
\usepackage{soul}
\usepackage{comment}

\usepackage[utf8]{inputenc}  
\usepackage[T1]{fontenc}  

\usepackage{booktabs} 

\usepackage{MnSymbol}

\usepackage{longtable}

\usepackage[all]{xy}

\usepackage[english]{babel} 
\usepackage[left=3.5cm,right=3.5cm,top=3.5cm,bottom=3cm,headsep=0.7cm]{geometry}
\parindent = 10pt

\usepackage{xargs}
%\usepackage[pdftex,dvipsnames]{xcolor}

\usepackage{csquotes}
\usepackage{stmaryrd}
% 
\usepackage[colorinlistoftodos,prependcaption,textsize=tiny]{todonotes}
\newcommandx{\unsure}[2][1=]{\todo[linecolor=red,backgroundcolor=red!25,bordercolor=red,#1]{#2}}
\newcommandx{\change}[2][1=]{\todo[linecolor=blue,backgroundcolor=blue!25,bordercolor=blue,#1]{#2}}
\newcommandx{\info}[2][1=]{\todo[linecolor=OliveGreen,backgroundcolor=OliveGreen!25,bordercolor=OliveGreen,#1]{#2}}
\newcommandx{\improvement}[2][1=]{\todo[linecolor=Plum,backgroundcolor=Plum!25,bordercolor=Plum,#1]{#2}}
\newcommandx{\thiswillnotshow}[2][1=]{\todo[disable,#1]{#2}}

 
%%%% Figures %%%%
 
\usepackage{pgf,tikz}
\usetikzlibrary{positioning, calc, arrows.meta}
\usetikzlibrary{arrows}
\usepackage{subcaption}
\usepackage{graphicx}
\usepackage[parfill]{parskip}

%%%% Enumeration %%%%

\usepackage{enumitem}

%%%% References %%%%

\usepackage{hyperref}
\usepackage[nameinlink]{cleveref}
\usepackage{mathtools}
\hypersetup{
    colorlinks,
    linkcolor={red!80!black},
    citecolor={blue!80!black}
}

\crefname{assumption}{assumption}{assumptions}  % Lowercase for mid-sentence usage
\Crefname{assumption}{Assumption}{Assumptions}  % Uppercase for starting sentences
\crefname{counterexample}{counterexample}{counterexamples}
\Crefname{counterexample}{Counterexample}{Counterexamples}

%%%% Environments %%%%

\theoremstyle{plain}
\begingroup
\theoremstyle{plain}
\newtheorem{theorem}{Theorem}[section]
\newtheorem{corollary}[theorem]{Corollary}
\newtheorem{proposition}[theorem]{Proposition}
\newtheorem{lemma}[theorem]{Lemma}
\theoremstyle{definition}
\newtheorem{definition}[theorem]{Definition}
\newtheorem{assumption}[theorem]{Assumption}
%\newtheorem{algo}[theorem]{Algorithm}
\theoremstyle{remark}
\newtheorem{remark}[theorem]{Remark}
\newtheorem{example}[theorem]{Example}
\newtheorem{counterexample}[theorem]{Counterexample}
\endgroup

\theoremstyle{definition}
\theoremstyle{remark}

\numberwithin{equation}{section}


%%%% Colors %%%%

\newcommand{\RRR}{\color{red}}
\newcommand{\EEE}{\color{black}}
\newcommand{\MMM}{\color{magenta}}
\newcommand{\blue}[1]{{\color{blue} #1}}

%%%% Symbols %%%%

% Special letters %
\newcommand{\R}{\mathbb{R}}
\newcommand{\N}{\mathbb{N}}
\newcommand{\NN}{\mathbb{N}}
\newcommand{\Z}{\mathbb{Z}}
\newcommand{\Q}{\mathbb{Q}} 

\newcommand{\Lip}{\mathrm{Lip}}
\newcommand{\MH}{\mathcal M_H}
\newcommand{\dH}{d_H}
\newcommand{\sinf}{\sigma_{\inf}}


% Corrected symbols %
\mathsurround=1pt
\mathchardef\emptyset="001F
\renewcommand{\d}{\mathrm{d}}
\newcommand{\de}{\partial}
\newcommand{\eps}{\varepsilon}
\renewcommand{\tilde}{\widetilde}
\renewcommand{\l}{\langle}
\renewcommand{\r}{\rangle}
\newcommand{\dist}{{\rm dist}}
\newcommand{\argmin}{\mathrm{argmin}}
\newcommand{\Tr}{\mathrm{Tr}}
\newcommand{\Lap}{\Delta}
\newcommand{\loc}{\mathrm{loc}}
\renewcommand{\b}{\boldsymbol{b}} 
\newcommand{\x}{\boldsymbol{x}} 
\newcommand{\rrho}{\boldsymbol{\rho}} 
\renewcommand{\u}{{\boldsymbol{u}}}
\newcommand{\relu}{\boldsymbol{\sigma}}
\newcommand{\K}{\mathcal{K}}
\newcommand{\X}{\mathcal{X}}
\newcommand{\Y}{\mathcal{Y}}
\newcommand{\F}{\mathcal{F}}
\newcommand{\D}{\mathcal{D}}
\newcommand{\C}{\mathbb{C}}
\newcommand{\A}{\mathcal{A}} 
\newcommand{\W}{\mathsf{W}} 
\newcommand{\E}{\mathbb{E}} 
\newcommand{\pphi}{\boldsymbol{\phi}} 
\newcommand{\ppsi}{\boldsymbol{\psi}} 
\renewcommand{\L}{\mathcal{L}}
\newcommand{\emp}{\hat{\pi}_N}
\newcommand{\RM}{\mathcal{RM}}
\newcommand{\Reg}{\mathrm{Reg}}

\newcommand{\Tree}{\mathsf{Tree}}
\newcommand{\TrTwo}{\Tree_2}
\newcommand{\Homeo}{\mathrm{Homeo}}
\newcommand{\sap}{\sigma_{\mathrm{ap}}}
\newcommand{\spp}{\sigma_{\mathrm{p}}}
\newcommand{\T}{\mathbb T}
\newcommand{\Uroots}[1]{\mathbb U_{#1}} % n-th roots of unity
\newcommand{\restr}{\!\upharpoonright\!}
\newcommand{\At}{\mathrm{At}}
\newcommand{\Perr}{\mathcal P}
\newcommand{\Lp}{L^p(\mathcal X,\omega)}
\newcommand{\Lq}{L^q(\mathcal X,\omega)}
\newcommand{\ip}[2]{\left\langle #1,#2\right\rangle}

\newcommand{\Top}{\mathbf{Top}}
\newcommand{\LCH}{\mathbf{LCH}}
\newcommand{\CGHaus}{\mathbf{CGHaus}}
\newcommand{\Rstar}{\mathbb{R}^{\star}}

\newcommand{\Baire}{\N^\N}
\newcommand{\dom}{\mathrm{dom}}
\newcommand{\id}{\mathrm{id}}
\newcommand{\pair}{\langle\,\cdot,\cdot\,\rangle}

\newcommand{\Bor}{\mathsf{Bor}}
\newcommand{\Cont}{\mathsf{Cont}}
\newcommand{\Comp}{\mathsf{Comp}}
\newcommand{\BaireMeas}{\mathsf{Baire}}
\newcommand{\SCIAbss}{\mathrm{SCI}^{\mathrm{BSS}}_A}

\newcommand{\pairWhite}[2]{\langle #1,#2\rangle}

\newcommand{\Lim}{\operatorname{lim}}
\newcommand{\Limi}[1]{\Lim^{(#1)}}

\newcommand{\SCIG}{\mathrm{SCI}_G}
\newcommand{\SCIA}{\mathrm{SCI}_A}

% Custom commands %
\newcommand{\vc}[2]{\begin{pmatrix}
    #1 \\ #2 
\end{pmatrix} }
\newcommand{\supp}{\mathrm{supp}}

\newcommand{\vct}[2]{\big( \begin{smallmatrix}
    #1 \\ #2 
\end{smallmatrix} \big) }

\newcommand{\mred}[1]{
\marginnote{\textcolor{red}{#1}}}

\def\eg{{\em e.g.}}
\def\cf{{\em cf.}}

\begin{document}

\title[Foundational Analysis Of The Solvability Complexity Index]{Foundational Analysis Of The Solvability Complexity Index: The Weihrauch-SCI Intermediate Hierarchy}
\titlecomment{An extended version is available as arXiv:2603.18955. It contains additional material on evaluation-induced topologies, Kihara-style piecewise-continuity reducibilities, and feasible real RAM models and their connection to the SCI}
\thanks{}	%optional

\author[C.~Sorg]{Christopher Sorg\lmcsorcid{0009-0003-9684-6966}}[a]
\address{University of the Bundeswehr Munich, Werner-Heisenberg-Weg 39, 85577 Neubiberg, Germany}	%optional
\email{chr.sorg@unibw.de}  %optional

\begin{abstract}
\noindent The Solvability Complexity Index (SCI) provides an extensional limit-height formalism for recovering a target map \(\Xi\) from finite samples of an evaluation interface \(\Lambda\subseteq\mathbb C^\Omega\) by finite-height towers of pointwise limits. We first give a foundational analysis of what this extensional framework does and does not determine.
We show that the SCI separation axiom is equivalent to a factorization of $\Xi$ through the full evaluation table, and we isolate the minimal logical role of $\Lambda$ as an information interface.

To connect the SCI to Type-2 computability and Weihrauch reducibility, we give an effective enrichment for countable $\Lambda$ by viewing the evaluation table image $I_{\Lambda}\subseteq\mathbb{C}^{\mathbb{N}}$ as a represented space and factoring $\Xi$ as $\widehat{\Xi}$.
We then define the Weihrauch-SCI rank of a problem as the least number of iterated limit-oracles needed to compute it in the Weihrauch sense, i.e.\ the least $k$ such that
$\widehat{\Xi}\le_{W}\lim^{(k)}$, and prove well-posedness and representation invariance of this rank.

A central negative result is that the unrestricted raw type-G SCI model (arbitrary post-processing of finite oracle transcripts) is generally not a computability model in the Type-2/Weihrauch sense: finite-query factorizations collapse raw type-G height, and analytic non-Borel decision problems yield examples with raw \(\SCIG=0\) but infinite Weihrauch-SCI rank. We therefore distinguish the raw extensional SCI from implemented SCI variants, where the indexed approximation table is required to be realized uniformly by a chosen class of operations.
To recover a robust bridge, we introduce an intermediate SCI hierarchy by restricting the admissible deepest-level post-processing to regularity classes (continuous/Borel/Baire) and, optionally, to fixed-query versus adaptive-query policies. We prove that these restrictions form hierarchies, and we establish comparison theorems showing what each restriction logically enforces.

Finally, we give self-contained canonical source problems over Cantor-matrix inputs which realize arbitrary finite standard raw type-G SCI heights. These examples are
not presented as computability models by themselves; rather, they are calibration objects for the extensional SCI and for the interaction between finite-query information, Borel hierarchy level, restricted SCI towers, and Weihrauch-style uniform iterated-limit complexity.
\end{abstract}

\maketitle

\section{Introduction}

The \textit{Solvability Complexity Index} (SCI) was introduced as a classification tool for
mathematical tasks that are not naturally captured by a single-shot notion of computation,
but instead admit (or require) \textit{towers of algorithms}: procedures whose outputs converge
to the desired value only after one or several limiting operations \cite{hansen2011solvability,ben2015computing,ColbrookHansen22}.
This perspective has proved particularly fruitful in computational spectral theory, where
the existence (or impossibility) of algorithms can depend delicately on the permitted input
interface and on the allowed post-processing of oracle information \cite{hansen2011solvability,ColbrookHansen22}.

At the same time, the SCI is intentionally \textit{extensional}: at its core, one fixes an input
class $\Omega$, a target map $\Xi:\Omega\to \mathcal{M}$ for a metric space $\mathcal{M}$, and an evaluation family $\Lambda$ (an
information interface), and defines what it means for a tower to compute $\Xi$ from finitely
many evaluation values (possibly iterated through limits). Here, extensional is not the set-theoretic axiom of extensionality,
the corresponding SCI-definitions just don't fix (on purpose) any syntax-level notation of computation yet.
This abstraction is a strength: it allows one to compare disparate numerical and information
models under a single roof but it also raises a foundational question that motivates the
present work:
\begin{align*}
\textbf{Q:} \text{ When does the SCI define a computability model?}
\end{align*}
A \enquote{computability model} is not merely a set of input-output tasks; it is a stable notion of
\textit{morphism} closed under composition and basic data operations, so that procedures can be
combined, relativized, and compared without ambiguity.
In Type-2 computable analysis this stability is enforced by \textit{representations} (names) and
\textit{realizers}: maps between represented spaces form a robust category in which computable
operations are automatically continuous at the level induced by the representations
\cite{weihrauch2000computable,brattka2021weihrauch}.
Weihrauch reducibility is meaningful precisely because it lives inside such a compositional
setting: it compares problems by uniform pre-/post-processing around a single oracle call,
and the resulting degree structures support a rich algebra of computational content
\cite{brattka2011weihrauch,brattka2021weihrauch}.

In contrast, the purely extensional SCI definition leaves open what counts as admissible
post-processing on finitely many oracle values.
Different interpretations of this interface can lead to different \enquote{heights} for the
same extensional task.
A classical example of such divergence is the gap between \textit{atomic} real-number interfaces
(which allow exact comparisons) and \textit{name-based} interfaces (Type-2), where computability
is necessarily continuous \cite{weihrauch2000computable,neumann2018topological,BSS98}.
This paper therefore takes the stance that statements of the form \enquote{$\Xi$ has SCI height $k$} 
should be read as \textit{conditional} on an explicit package of premises describing 
(i) the information interface, and (ii) the closure/regularity properties
demanded of the admissible post-processing.

A second distinction is equally important. A tower may be understood merely as an
extensional family of approximants, or as a uniformly implemented approximation
scheme.  In the first reading, each indexed approximant may be specified separately;
this is sufficient for an information-theoretic limit-height classification, but it is not
a computability model in the represented-space sense. In the second reading, the
whole indexed table of approximants must be produced uniformly by one admissible
procedure, for instance by one BSS machine, one Type-2 functional, or one realizer in
a chosen regularity class. Weihrauch comparability belongs to this second reading.

A tempting repair would be to require every intermediate semantic limit map in an
SCI tower to be a finite-information general algorithm. \Cref{prop:all-level-collapse} shows why this is
not the right solution: such an all-level repair is too strong and collapses the finite
positive hierarchy to height \(1\). The correct repair is instead uniform
implementation of the indexed approximation table.
There is an instructive parallel here to a familiar point in logic.
On the proof-theoretic side, theoremhood (\enquote{tautology}) is defined relative to a chosen
deductive calculus (axioms and inference rules); on the semantic side, validity is defined
relative to a chosen account of consequence (e.g.\ Tarski's model-theoretic criterion)
\cite{Tarski02,shapiro1998logical}.
Carnap famously emphasized that different formal frameworks may be adopted for different
purposes (\enquote{tolerance}), while Quine challenged sharp analytic/synthetic boundaries and
thus any overly rigid separation between \enquote{truth by meaning} and \enquote{truth by fact}
\cite{carnap2014logical,quine2000two}.
In the contemporary literature, these tensions also motivate more pluralist readings of
logical consequence \cite{beall2000logical}, as well as detailed debates about what Tarski's
definition does and does not capture \cite{etchemendy1990concept}.
The moral for the present context is not that the SCI is purely philosophical, but that \textit{the
informativeness of any classification depends on making its background commitments explicit}.
Our goal is to do this in a way that stays mathematically sharp while remaining neutral with
respect to foundational preferences.

\textbf{Overview of the paper.}$\,$
\Cref{sec:SCIdef} fixes the standard extensional SCI and separates it from uniformly
implemented tower readings. \Cref{sec:MinLogRoleEv} identifies the minimal logical role of the
evaluation family \(\Lambda\): the SCI consistency condition is equivalent to
factorization of the target through the full evaluation table. \Cref{sec:EffEnrich} enriches
countable evaluation interfaces by represented information spaces and defines the
Weihrauch-SCI rank. \Cref{sec:G-Case} proves that unrestricted raw type-\(G\) SCI is not
generally comparable with Weihrauch or Type-2 complexity: finite-query factorizations
collapse raw \(\SCIG\), while analytic non-Borel decision problems have infinite
Weihrauch--SCI rank. \Cref{sec:IntHierarchy} introduces the regularity-restricted intermediate SCI
hierarchy and proves the basic comparison theorems. \Cref{sec:ExampleSCI} gives canonical
Cantor-matrix source problems of exact finite raw type-\(G\) height and a tagged union
of infinite height. \Cref{sec:uniformity-minimality} collects the uniformity and minimality results: it explains
why all-level finite-information towers are too restrictive, why non-uniform
deepest-only towers are too weak for Weihrauch comparability, and why the pure
\(\mathcal R\)-\(\Lim\) normal form gives the correct represented-space bridge.

\section{The Standard Extensional SCI And Implemented Readings}{\label{sec:SCIdef}}

We are interested in measuring the \enquote{algorithmic complexity} of mathematical spectral approximation problems in the sense of the \textit{Solvability Complexity Index} (SCI), which measures the minimal number of nested limits required by any algorithm to solve a given (infinite-dimensional) problem \cite{hansen2011solvability,ben2015computing,ben2015new,ColbrookHansen22}. The SCI framework has revealed sharp boundaries for spectral computations of (unbounded) operators. Some spectral problems are provably not solvable by any single-limit procedure (hence admit no global error control), yet become solvable by towers of algorithms of finite height or may not be solvable for a finite height at all, in the latter the SCI being therefore infinity, see e.g.\ \cite[Theorem B.1]{colbrook2024limits}. We first state the standard extensional SCI. We then distinguish it from implemented variants, where the indexed approximation table is required to be produced uniformly by a specified computation model. This distinction is logically necessary for Weihrauch comparisons.

\textbf{Notation:} In this work, we will always use the notation $\NN := \{1,2,3,\dots\}$ and $\NN_0:=\NN \cup \{0\}$.

\begin{definition}[Metric space]\label{def:metric}
A \textit{metric space} is a pair $(\mathcal{M},d)$ where $\mathcal{M}$ is a set and the map
$d:\mathcal{M}\times\mathcal{M}\to [0,\infty)$ satisfies $\forall x,y,z \in \mathcal{M}$:
(i) $d(x,y)=0 \Leftrightarrow x=y$;
(ii) $d(x,y)=d(y,x)$;
(iii) $d(x,z)\le d(x,y)+d(y,z)$.
\end{definition}

\begin{definition}[Complex-valued evaluation set]\label{def:eval-set}
Let $\Omega$ be a set. Write $\C^\Omega:=\{f:\Omega\to\C\}$.
An \textit{evaluation set} on $\Omega$ is a subset $\Lambda\subseteq \C^\Omega$.
Equivalently, $\Lambda$ is a family of maps $f:\Omega\to\C$.
\end{definition}

\begin{remark}[Non-$\C$-valued \enquote{natural} evaluations]\label{rem:nonC}
Many evaluations are not $\C$-valued.
To fit \cref{def:eval-set} \textit{literally}, one fixes once and for all an injective encoding
$\iota:V\hookrightarrow \C$ of the value set $V$ (or into $\C^m$ and then into $\C$),
and replaces each $g:\Omega\to V$ by $\iota\circ g:\Omega\to\C$.
This does not change the \textit{information content} of the oracle family as long as $\iota$ is injective.
\end{remark}

The now following SCI definitions are based on the definitions in \cite[Appendix~A]{colbrook2024limits} but slightly precised for our purposes in this work.

\begin{definition}[Computational problem; consistency \enquote{axiom}]\label{def:comp-prob}
A \textit{computational problem} is a quadruple
\[
  \mathcal{P}=\bigl( \Xi,\Omega,(\mathcal{M},d),\Lambda \bigr),
\]
where
\begin{itemize}
\item $\Omega$ is a set (the \textit{primary set}, i.e.\ input class);
\item $(\mathcal{M},d)$ is a metric space in the sense of \cref{def:metric};
\item $\Xi:\Omega\to\mathcal{M}$ is a function (the \textit{problem function});
\item $\Lambda\subseteq \C^\Omega$ is an evaluation set in the sense of \cref{def:eval-set}.
\end{itemize}
We impose the \textit{consistency \enquote{axiom}}:
\begin{equation}\label{eq:consistency}
\forall A,B\in\Omega,\quad \Xi(A)\neq \Xi(B)\ \Longrightarrow\ \exists f\in\Lambda\ \text{s.t. } f(A)\neq f(B).
\end{equation}
\end{definition}

For any set $X$, we will write in the following
\[
[X]^{<\omega} := \{E \subseteq X : E \text{ is finite}\}.
\]

\begin{definition}[General algorithm]\label{def:gen-alg}
Let $\mathcal{P}=\bigl( \Xi,\Omega,(\mathcal{M},d),\Lambda \bigr)$ be a computational problem.
A \textit{general algorithm} for $\mathcal{P}$ is a pair $(\Gamma,\Lambda_\Gamma)$ where
\[
\Gamma:\Omega\to \mathcal{M}
\quad\text{and}\quad
\Lambda_\Gamma:\Omega\to [\Lambda]^{<\omega}
\]
assigns to each $A\in\Omega$ a finite nonempty set of queries $\Lambda_\Gamma(A)\subseteq\Lambda$
such that for all $A,B\in\Omega$:\hfill
\begin{enumerate}
\item (\textbf{Finite-information dependence}): If $f(B)=f(A)$ for all $f\in\Lambda_\Gamma(A)$,
then $\Gamma(B)=\Gamma(A)$.
\item (\textbf{Stability of the query set}): If $f(B)=f(A)$ for all $f\in\Lambda_\Gamma(A)$,
then $\Lambda_\Gamma(B)=\Lambda_\Gamma(A)$.
\end{enumerate}
\end{definition}

\begin{definition}[Query policy; fixed vs.\ adaptive querying]\label{def:query-policy}
Let $\mathcal{P}=\bigl( \Xi,\Omega,(\mathcal{M},d),\Lambda \bigr)$ be an SCI computational problem and let $(\Gamma,\Lambda_\Gamma)$
be a general algorithm for $\mathcal{P}$ in the sense of \cref{def:gen-alg}.

The map $\Lambda_\Gamma:\Omega\to[\Lambda]^{<\omega}$ is called the
\textit{query policy} of the algorithm.

\begin{enumerate}
\item We call $(\Gamma,\Lambda_\Gamma)$ \textit{fixed-query} if there exists a finite set
$Q\subseteq\Lambda$ such that $\Lambda_\Gamma(A)=Q$ for all $A\in\Omega$.
\item Otherwise we call $(\Gamma,\Lambda_\Gamma)$ \textit{adaptive-query}.
Equivalently, it is adaptive-query iff there exist $A,B\in\Omega$ with
$\Lambda_\Gamma(A)\neq\Lambda_\Gamma(B)$.
\end{enumerate}
\end{definition}

\begin{remark}[Extensional meaning of adaptivity]\label{rem:extensional-adaptivity}
\Cref{def:gen-alg} is the extensional footprint of adaptive oracle access: the set of
queries used on input $A$ is determined solely by the answers to \textit{those} queries.
Formally, for all $A,B\in\Omega$,
\[
\Bigl(\forall f\in\Lambda_\Gamma(A)\;\; f(B)=f(A)\Bigr)
\ \Longrightarrow\
\Bigl(\Lambda_\Gamma(B)=\Lambda_\Gamma(A)\ \wedge\ \Gamma(B)=\Gamma(A)\Bigr).
\]
Thus two inputs that induce the same query-answer transcript on the queried evaluations
force identical algorithmic behaviour (same queries, same output).
\end{remark}

\begin{remark}[Suppressing the query policy]\label{rem:suppress-query-policy}
When only the output map matters, we will (ab)use notation and write $\Gamma$ for the
pair $(\Gamma,\Lambda_\Gamma)$, keeping $\Lambda_\Gamma$ implicit; when needed we
write $\Lambda_\Gamma(A)$ explicitly.
This avoids ambiguity later when towers are written only in terms of their output maps.
\end{remark}

\begin{definition}[Standard extensional tower of algorithms]\label{def:standard-tower}
Let
\[
        \mathcal P=(\Xi,\Omega,(\mathcal M,d),\Lambda)
\]
be an SCI computational problem and let \(k\in\mathbb N_0\).

A standard tower of height \(0\) for \(\mathcal P\) is a single general algorithm $(\Gamma,\Lambda_\Gamma)$ such that
\[
        \Gamma(A)=\Xi(A)
\]
for $A\in\Omega$. For \(k\in \mathbb{N}\), a standard tower of height \(k\) for \(\mathcal P\) consists of general algorithms
\[
        \Gamma_{n_k,\ldots,n_1}:\Omega\to \mathcal M, \, (n_1,\ldots,n_k)\in\mathbb N^k,
\]
such that for every \(A\in\Omega\) the iterated limit
\[
        \Xi(A) = \lim_{n_k\to\infty} \lim_{n_{k-1}\to\infty} \cdots \lim_{n_1\to\infty} \Gamma_{n_k,\ldots,n_1}(A)
\]
exists in the metric \(d\).

Equivalently, the intermediate maps are defined recursively, whenever the displayed limits exist, by
\[
        \Gamma_{n_k,\ldots,n_r}(A) := \lim_{n_{r-1}\to\infty} \Gamma_{n_k,\ldots,n_{r-1}}(A),
\]
where $r=2,\ldots,k$, and
\[
        \Gamma_{n_k}(A) := \lim_{n_{k-1}\to\infty} \Gamma_{n_k,n_{k-1}}(A).
\]
These intermediate maps are semantic pointwise limits. They are not required to be general algorithms. Only the deepest-level maps
\[
        \Gamma_{n_k,\ldots,n_1}
\]
are required to be general algorithms.
\end{definition}

\begin{definition}[Uniformly implemented tower]\label{def:uniformly-implemented-tower}
Let
\[
        \mathcal P=(\Xi,\Omega,(\mathcal M,d),\Lambda)
\]
be an SCI computational problem, and let \(\mathfrak C\) be a specified implementation class.

For \(k=0\), a height-\(0\) \(\mathfrak C\)-implemented tower for \(\mathcal P\) is a single \(\mathfrak C\)-admissible map
\[
        \mathcal A:\Omega\to \mathcal M
\]
such that
\[
        \mathcal A(A)=\Xi(A)
\]
for $A\in\Omega$.

For \(k\ge1\), a height-\(k\) \(\mathfrak C\)-implemented tower for \(\mathcal P\) consists of a single \(\mathfrak C\)-admissible indexed approximation scheme
\[
        \mathcal A:\mathbb N^k\times\Omega\to \mathcal M, \, (n_k,\ldots,n_1,A)\mapsto
        \mathcal A(n_k,\ldots,n_1;A),
\]
such that for every \(A\in\Omega\),
\[
        \Xi(A) =
        \lim_{n_k\to\infty}
        \lim_{n_{k-1}\to\infty}
        \cdots
        \lim_{n_1\to\infty}
        \mathcal A(n_k,\ldots,n_1;A).
\]

The phrase \(\mathfrak C\)-admissible means that the whole indexed table is implemented
uniformly in the indices and in the input according to the model \(\mathfrak C\).
For example:
\begin{enumerate}
\item for a BSS implementation, \(\mathcal A\) is produced by one oracle BSS machine
which receives \((n_k,\ldots,n_1)\) as ordinary input and has oracle access to \(A\);
\item for a Type-2 implementation, \(\mathcal A\) is realized by one Type-2 functional
which receives a name of \(A\) and the indices and outputs a name of the approximant;
\item for a regularity-restricted implementation, the indexed approximation map is
required to have the prescribed regularity uniformly in the indices.
\end{enumerate}
\end{definition}

Now we state first the tower definitions as in most of the SCI-literature.
\begin{definition}[Historical informal tower types in the SCI literature]\label{def:tower-types}
One may impose restrictions on the deepest-level approximants in a standard tower, or one may impose a uniform implementation requirement on the whole indexed table.
\begin{itemize}
\item \textbf{General towers ($\alpha=G$):} no restriction beyond \cref{def:standard-tower}.
\item \textbf{Arithmetic towers ($\alpha=A$):} A map $\Gamma:\Omega\to\mathcal{M}$ is called an \textit{arithmetic algorithm} if
$\Gamma$ is a general algorithm in the sense of \cref{def:gen-alg} and, moreover,
for every $A\in\Omega$ the value $\Gamma(A)$ is obtained from the finite tuple
$\bigl(f(A)\bigr)_{f\in\Lambda_\Gamma(A)}$ by a finite computation using only
\textit{arithmetic operations} on $\C$. We call a tower of algorithms in the sense of \cref{def:standard-tower} an \textit{arithmetic tower}, 
if each deepest-level map $\Gamma_{n_k,\ldots,n_1}$ is an arithmetic algorithm.
\end{itemize}
\end{definition}

We note that for the purposes of this work \cref{def:tower-types} is \textit{not} precise enough. First \enquote{finite computation} is not a mathematical object unless one chooses a machine model (or equivalent syntax); \enquote{comparisons} over $\C$ are underspecified and for general $\mathcal M$ it is unclear what it means to \enquote{output} an element of $\mathcal M$. The most natural fix we could think of is a BSS-specification, which we will make precise now. This precision is in line with the more precise SCI-definitions provided e.g.\ in \cite[\S 2]{colbrook2020foundations} or in \cite[\S 6]{ben2015computing}.

\begin{definition}[$\R^\star$ (finite real tuples)]\label{def:Rstar}
Let
\[
\R^\star \ :=\ \bigsqcup_{m\in\N} \R^m
\]
be the disjoint union of all finite-dimensional Euclidean spaces.
\end{definition}

\begin{definition}[Finite output interface]\label{def:finite-output-interface}
Let $(\mathcal{M},d)$ be a metric space.
A \textit{finite output interface} for $\mathcal{M}$ is a fixed (possibly partial) map
\[
\delta^{\mathrm{fin}}_{\mathcal{M}}:\subseteq \R^\star \to \mathcal{M}.
\]
Elements $u\in\dom(\delta^{\mathrm{fin}}_{\mathcal{M}})$ are called \textit{finite codes} for
$\delta^{\mathrm{fin}}_{\mathcal{M}}(u)\in\mathcal{M}$.
\end{definition}

\begin{definition}[Evaluation oracle for a countable evaluation set]\label{def:eval-oracle}
Assume $\Lambda$ is countable and fix an enumeration $\Lambda=(f_i)_{i\in\N}$.
For each $A\in\Omega$ define the \textit{evaluation oracle}
\[
\mathcal{O}_A:\N\to\R^2,\qquad \mathcal{O}_A(i):=\bigl(\operatorname{Re}(f_i(A)),\operatorname{Im}(f_i(A))\bigr).
\]
\end{definition}

\begin{definition}[Oracle BSS machine (parameter-free; partial QUERY)]\label{def:oracle-bss}
Fix an enumeration $\Lambda=(f_i)_{i\in\N}$ and the associated oracle maps
$\mathcal{O}_A:\N\to\R^2$ from \cref{def:eval-oracle}.
An \textit{oracle BSS machine over $(\R,+,-,\cdot,/,<)$} in the sense of \cite{BSS98}
\textit{without real parameters} is a BSS machine whose only built-in constants are
$0$ and $1$ (hence all rational constants are computable inside the machine),
equipped with an additional instruction
\[
\mathrm{QUERY}: \
\text{read } i \text{ from a designated register and write } \mathcal{O}_A(i)\in\R^2
\text{ into designated registers.}
\]
The instruction $\mathrm{QUERY}$ is \textit{partial}: it is defined only if the designated
register contains a natural number $i\in\N$ (encoded as the real $i$); otherwise the run
is undefined (equivalently: the machine diverges).

A run \textit{halts} if it reaches a halting state after finitely many steps; upon halting it
outputs a finite tuple $u\in\R^\star$ (encoded in the standard BSS way).
\end{definition}

This is the formal replacement for \enquote{finite computation using arithmetic operations and comparisons}.
\begin{definition}[BSS-arithmetic algorithm]\label{def:BSS-arith-alg}
Let $\mathcal{P}=\bigl( \Xi,\Omega,(\mathcal{M},d),\Lambda \bigr)$ be a computational problem.
Assume:
\begin{enumerate}
\item $\Lambda$ is countable with a fixed enumeration $\Lambda=(f_i)_{i\in\N}$;
\item $\mathcal{M}$ is equipped with a fixed finite output interface
      $\delta^{\mathrm{fin}}_{\mathcal{M}}$ as in \cref{def:finite-output-interface}.
\end{enumerate}
A map $\Gamma:\Omega\to\mathcal{M}$ is called a \textit{BSS-arithmetic algorithm}
if there exists an oracle BSS machine $\mathbb{M}$ in the sense of \cref{def:oracle-bss}
such that for every $A\in\Omega$ the run of $\mathbb{M}$ with oracle $\mathcal{O}_A$
halts and outputs some $u\in\dom(\delta^{\mathrm{fin}}_{\mathcal{M}})$ with
\[
\delta^{\mathrm{fin}}_{\mathcal{M}}(u)=\Gamma(A).
\]
We require $\mathbb M$ to be parameter-free in the sense of \cref{def:oracle-bss}.
\end{definition}

Now we can give a formal precise version of \cref{def:tower-types}.
\begin{definition}[Tower types: general (G) vs.\ arithmetic (A)]\label{def:tower-types-BSS}
Let
\[
        \mathcal P=(\Xi,\Omega,(\mathcal M,d),\Lambda)
\]
be an SCI computational problem.
\begin{itemize}
\item \textbf{Raw general towers ($\alpha=G$):} A type-G tower is a standard tower in the sense of \cref{def:standard-tower}.
\item \textbf{Uniformly BSS-arithmetic towers ($\alpha=A$):} assume $\Lambda$ is countable with fixed enumeration and
$(\mathcal{M},d)$ comes with a fixed finite output interface as in
\cref{def:BSS-arith-alg}. A height-\(k\) BSS-arithmetic tower is a standard tower whose deepest-level approximants
\[
        \Gamma_{n_k,\ldots,n_1}
\]
are produced uniformly by one parameter-free oracle BSS machine. More precisely, there is a single oracle BSS machine \(\mathbb M\) such that, for every
\[
        (n_k,\ldots,n_1)\in\mathbb N^k
        \quad\text{and every }A\in\Omega,
\]
the machine \(\mathbb M\), given the indices \((n_k,\ldots,n_1)\) in its input
registers and oracle access to \(O_A\), halts and outputs a code
\[
        u_{n_k,\ldots,n_1}(A)\in\operatorname{dom}(\delta_\mathcal M^{\rm fin})
\]
with
\[
        \delta_\mathcal M^{\rm fin}(u_{n_k,\ldots,n_1}(A))
        =
        \Gamma_{n_k,\ldots,n_1}(A).
\]
For height \(0\), the same definition is read with no index input.
\end{itemize}
\end{definition}

\begin{definition}[SCI]\label{def:sci}
Let \(\mathcal P\) be an SCI computational problem.

\begin{enumerate}
\item[(S1)] The raw type-G SCI is defined by
\[
        \SCIG(\mathcal P)
        :=
        \min\{k\in\mathbb N_0:
        \text{there exists a standard raw type-G tower of height }k\text{ for } \mathcal P\},
\]
with value \(\infty\) if the set is empty.

\item[(S2)] When the BSS-arithmetic implementation of \cref{def:tower-types-BSS} is available, define
\[
        \SCIA^{\rm BSS}(\mathcal P)
        :=
        \min\{k\in\mathbb N_0:
        \text{there exists a uniformly BSS-arithmetic tower of height }k\text{ for } \mathcal P\},
\]
with value \(\infty\) if the set is empty.
\end{enumerate}
\end{definition}

We write $\SCIAbss(\mathcal P)$ for $\mathrm{SCI}_A(\mathcal P)$ when the tower type $\alpha=A$
is interpreted via BSS-arithmetic base maps as in \cref{def:BSS-arith-alg,def:tower-types-BSS}.

\begin{remark}[Why extra structure and uniformity are logically necessary]\label{rem:need-extra}
\Cref{def:comp-prob,def:gen-alg,def:standard-tower} are \textit{purely extensional} in that they
are statements about maps between underlying sets. There are two distinct missing ingredients for a represented-space computability model. First, unrestricted type-G post-processing allows arbitrary maps on finite transcripts; therefore finite-query factorizations collapse raw \(\SCIG\) to height \(0\). Second, even if the deepest-level approximants are individually computable, one still
needs uniformity in the tower indices. A computability model should provide one procedure producing the indexed approximation table, not a separately chosen procedure for every multi-index. This is the role of implemented towers in the sense of \cref{def:uniformly-implemented-tower} and, in the Weihrauch setting, of the pure $\mathcal R-\Lim$ normal form in \cref{sec:uniformity-minimality}.
\end{remark}

\begin{definition}[All-level finite-information tower]\label{def:all-level-tower}
An all-level finite-information tower is a standard tower in which every intermediate semantic limit map
\[
        \Gamma_{n_k},\Gamma_{n_k,n_{k-1}},\ldots,\Gamma_{n_k,\ldots,n_2}
\]
is also required to be a general algorithm in the sense of \cref{def:gen-alg}.
\end{definition}

\begin{proposition}[Collapse of finite all-level towers]\label{prop:all-level-collapse}
Let \(\mathrm{SCI}^{\rm all}_G(\mathcal P)\) denote the least height of an all-level finite-information tower computing \(\mathcal P\). Then, for every SCI problem \(\mathcal P\),
\[
        0<\mathrm{SCI}^{\rm all}_G(\mathcal P)<\infty \quad\Longrightarrow\quad \mathrm{SCI}^{\rm all}_G(\mathcal P)=1.
\]
Equivalently, the all-level finite hierarchy has only the finite values \(0\) and \(1\).
\end{proposition}

\begin{proof}
Suppose \(\mathcal P\) has an all-level tower of height \(k\in \mathbb{N}\). By definition, the outermost maps
\[
        \Gamma_{n_k}:\Omega\to \mathcal M
\]
are themselves general algorithms, and
\[
        \Xi(A)=\lim_{n_k\to\infty}\Gamma_{n_k}(A)
\]
for $A\in\Omega$. Thus \((\Gamma_{n_k})_{n_k\in\mathbb N}\) is already a height-\(1\) type-G tower. Hence
\[
        \mathrm{SCI}^{\rm all}_G( \mathcal P)\le1.
\]
If the height is positive, it is not \(0\), so it must be \(1\).
\end{proof}

\section{The Minimal Logical Role Of The Evaluation Set \texorpdfstring{$\Lambda$}{Lambda}}{\label{sec:MinLogRoleEv}}

\begin{definition}[Full evaluation table and $\Lambda$-indistinguishability]
For each $f\in\Lambda$ let $D_f$ denote its codomain. Define the \textit{full evaluation table map}
\[
  \mathrm{Ev}_\Lambda:\Omega \to \prod_{f\in\Lambda} D_f,
  \qquad
  \mathrm{Ev}_\Lambda(A) := (f(A))_{f\in\Lambda}.
\]
(If all $D_f=\mathbb{C}$, then $\prod_{f\in\Lambda}D_f = \mathbb{C}^\Lambda$.)

Define an equivalence relation $\sim_\Lambda$ on $\Omega$ by
\[
  A\sim_\Lambda B \quad:\Longleftrightarrow\quad \forall f\in\Lambda,\ f(A)=f(B).
\]
\end{definition}

\begin{theorem}[Consistency $\Leftrightarrow$ factorization through the evaluation table]\label{th:consEvTable}
Let $\mathcal{P}=\bigl( \Xi,\Omega,(\mathcal{M},d),\Lambda \bigr)$.
The following are equivalent:
\begin{enumerate}
  \item (Consistency) $\Xi(A)\neq\Xi(B)\Rightarrow \exists f\in\Lambda:\ f(A)\neq f(B)$.
  \item (Extensionality) $A\sim_\Lambda B \Rightarrow \Xi(A)=\Xi(B)$.
  \item (Factorization) There exists a unique map
    \[
      \Phi:\mathrm{Ev}_\Lambda(\Omega)\to \mathcal{M}
    \]
    such that $\Xi=\Phi\circ \mathrm{Ev}_\Lambda$.
\end{enumerate}
\end{theorem}

\begin{proof}
(1)$\Leftrightarrow$(2) is contrapositive.
Assume (2). Define $\Phi$ on $\mathrm{Ev}_\Lambda(\Omega)$ by
$\Phi(\mathrm{Ev}_\Lambda(A)):=\Xi(A)$.
This is well-defined: if $\mathrm{Ev}_\Lambda(A)=\mathrm{Ev}_\Lambda(B)$ then $A\sim_\Lambda B$, hence $\Xi(A)=\Xi(B)$.
Then $\Xi=\Phi\circ \mathrm{Ev}_\Lambda$ by construction.
Uniqueness holds because $\mathrm{Ev}_\Lambda(\Omega)$ is the image.
Conversely, (3) implies (2) because equality of evaluation tables forces equality of $\Phi$-images.
\end{proof}

\begin{remark}[Evaluation-induced topology]
An evaluation family \(\Lambda\) induces the initial topology generated by the maps in
\(\Lambda\). For point evaluations on \(C(X,Y)\), this is the topology of pointwise
convergence, not the compact-open topology. Thus finite point-evaluation information
does not automatically provide the uniform control usually associated with compact-open
or compactly generated function-space topologies.
\end{remark}


\begin{corollary}[Necessity of consistency]\label{cor:NecConsis}
If consistency fails (i.e.\ there exist $A,B\in\Omega$ with $A\sim_\Lambda B$ but $\Xi(A)\neq\Xi(B)$),
then there is no tower of general algorithms (of any finite height) computing $\Xi$ from $\Lambda$.
Hence $\mathrm{SCI}_G(\mathcal{P})=\infty$.
\end{corollary}

\begin{proof}
Let $\{\Gamma_{n_k,\dots,n_1}\}$ be the deepest-level maps of any proposed finite-height tower.
Each deepest-level general algorithm $\Gamma_{n_k,\dots,n_1}$ depends only on finitely many evaluations, so if $A\sim_\Lambda B$ then in particular
$f(A)=f(B)$ for all $f\in\Lambda_{\Gamma_{n_k,\dots,n_1}}(A)$.
By \cref{def:gen-alg}(1), $\Gamma_{n_k,\dots,n_1}(A)=\Gamma_{n_k,\dots,n_1}(B)$
for all indices, hence every iterated limit agrees at $A$ and $B$, contradicting $\Xi(A)\neq\Xi(B)$.
\end{proof}

\section{Effective Enrichment: Represented Spaces Induced By Countable \texorpdfstring{$\Lambda$}{Lambda}}{\label{sec:EffEnrich}}

We recall some basic definitions from Weihrauch theory.

\begin{definition}[Represented space]\label{def:rep-space}
A \textit{represented space} is a pair $X=(|X|,\delta_X)$ where
$\delta_X:\subseteq \Baire \to |X|$ is a partial surjection.
A $p\in\Baire$ with $\delta_X(p)=x$ is called a \textit{name} of $x$.
\end{definition}

\begin{definition}[Standard representations for $\C$ and countable products]\label{def:rep-C-product}
Fix a computable dense sequence
\[
Q:\mathbb{N}\to\mathbb{C}.
\]
Let $\delta_{\mathbb{C}}$ be the standard Cauchy representation of $\mathbb{C}$ induced by
the usual metric and the dense sequence $Q$.
(Equivalently, any Cauchy representation induced by the usual metric and a computable
dense sequence; any two such Cauchy representations are computably isomorphic.)

For a sequence of represented spaces $(D_n,\delta_n)_{n\in\N}$, define the product
representation $\delta_{\prod D_n}$ on $\prod_n D_n$ via any fixed computable coding
of sequences of names into one name. In particular, this yields the standard product
representation $\delta_{\C^\N}$ on $\C^\N$.
\end{definition}

\begin{definition}[Countable evaluation set and evaluation-table map]\label{def:countable-Lambda}
Let $\Lambda\subseteq\C^\Omega$ be \textit{countable}. Fix once and for all
\[
\Lambda=\{f_n\}_{n\in\N},\qquad f_n:\Omega\to\C.
\]
Define the \textit{evaluation-table map}
\[
\mathrm{Ev}_\Lambda:\Omega\to \C^\N,\qquad \mathrm{Ev}_\Lambda(A):=(f_n(A))_{n\in\N}.
\]
Define the \textit{$\Lambda$-information set}
\[
\mathcal{I}_\Lambda:=\mathrm{Ev}_\Lambda(\Omega)\subseteq \C^\N.
\]
Equip $\C^\N$ with the product representation $\delta_{\C^\N}$ in the sense of \cref{def:rep-C-product},
and equip $\mathcal{I}_\Lambda$ with the induced subspace representation, denoted $\delta_{\mathcal{I}_\Lambda}$.
\end{definition}

\begin{lemma}[Consistency $\Rightarrow$ $\Xi$ factors through $\mathcal{I}_\Lambda$]\label{lem:factor}
Assume the consistency \enquote{axiom} holds for $\mathcal{P}=\bigl( \Xi,\Omega,(\mathcal{M},d),\Lambda \bigr)$ and that
$\Lambda=\{f_n\}_{n\in\mathbb{N}}$ is countable. Let $\mathrm{Ev}_\Lambda:\Omega\to\mathbb{C}^{\mathbb{N}}$
be the evaluation-table map and $I_\Lambda:=\mathrm{Ev}_\Lambda(\Omega)$.
Then there is a unique map $\widehat{\Xi}:I_\Lambda\to \mathcal{M}$ such that
$\widehat{\Xi}(\mathrm{Ev}_\Lambda(A))=\Xi(A)$ for all $A\in\Omega$.
\end{lemma}

\begin{proof}
This is up to notation the exact same statement as \cref{th:consEvTable}(3), still we give again a proof here for completeness. Define $\widehat{\Xi}:I_\Lambda\to \mathcal{M}$ by
\[
\widehat{\Xi}(x):=\Xi(A)\quad\text{for any }A\in\Omega\text{ with }\mathrm{Ev}_\Lambda(A)=x.
\]
We must show this is well-defined. Suppose $A,B\in\Omega$ satisfy
$\mathrm{Ev}_\Lambda(A)=\mathrm{Ev}_\Lambda(B)$. Then $f_n(A)=f_n(B)$ for every $n$,
hence $f(A)=f(B)$ for every $f\in\Lambda$. By the consistency \enquote{axiom}, this implies
$\Xi(A)=\Xi(B)$. Therefore $\widehat{\Xi}$ is well-defined.

By construction, $\widehat{\Xi}(\mathrm{Ev}_\Lambda(A))=\Xi(A)$ for all $A$.
Uniqueness follows because $\mathrm{Ev}_\Lambda:\Omega\to I_\Lambda$ is surjective.
\end{proof}

\begin{remark}[Enumeration-(in)dependence]\label{rem:enum-indep}
Let $(f_n)_{n\in\N}$ and $(g_n)_{n\in\N}$ be \textit{bijections}
$\mathbb{N}\to\Lambda$ (i.e.\ enumerations without repetition) of the same countable set $\Lambda$.
Then there exists a (unique) permutation $\pi:\N\to\N$ such that $g_n=f_{\pi(n)}$ for all $n$.

The induced coordinate permutation
\[
P_\pi:\C^\N\to\C^\N,\qquad P_\pi\bigl((z_n)_n\bigr):=(z_{\pi(n)})_n
\]
is always a homeomorphism (for the product topology). Moreover, with respect to the
standard product representation on $\C^\N$, the map $P_\pi$ is computable (and hence
a computable isomorphism of represented spaces) iff $\pi$ is a computable
function.

Consequently, $\mathcal{I}_\Lambda\subseteq\C^\N$ obtained from two enumerations are
always homeomorphic, but they are computably isomorphic only under computable reindexings.
\end{remark}

\subsection{Weihrauch Reducibility, Iterated Limits, And The SCI-Rank}

\begin{definition}[Pairing/tupling function]\label{def:tupling}
Fix once and for all a computable bijection
\[
\langle\cdot\rangle_k:\N^{k+1}\to\N
\]
for each $k\in\N$ (coding $(n_1,\dots,n_k,m)$ into one natural number).
\end{definition}

\begin{definition}[The limit operator $\Lim$ on Baire space]\label{def:lim}
Fix a computable bijection $\langle\cdot,\cdot\rangle:\NN^2\to\NN$.
For $p\in\NN^\NN$ write $p_n(k):=p(\langle n,k\rangle)$.
Define $\Lim:\subseteq\NN^\NN\to\NN^\NN$ by
\[
p\in\dom(\Lim)\;:\Longleftrightarrow\;(\forall k\in\NN)\; \text{the sequence }n\mapsto p_n(k)
\text{ is eventually constant},
\]
and then
\[
(\Lim(p))(k):=\lim_{n\to\infty} p_n(k).
\]
\end{definition}

\begin{definition}[Iterates of $\Lim$]\label{def:lim-iterates}
Define $\Limi{0}:=\id_{\NN^\NN}$ and $\Limi{n+1}:=\Lim\circ\Limi{n}$ for $n\in\NN_0$.
\end{definition}

\begin{definition}[Realizer]
Let $f:\subseteq X \rightrightarrows Y$ be a multi-valued map between represented spaces
$X=(|X|,\delta_X)$ and $Y=(|Y|,\delta_Y)$. A partial map
\[
F:\subseteq \mathbb{N}^{\mathbb{N}} \to \mathbb{N}^{\mathbb{N}}
\]
realizes $f$, written $F \vdash f$, if for every $p \in \dom(\delta_X)$ with
$\delta_X(p)\in\dom(f)$ one has $\delta_Y(F(p)) \in f(\delta_X(p))$.
\end{definition}

\begin{definition}[Weihrauch reducibility on represented spaces]\label{def:weihrauch}
Let $f:\subseteq X\rightrightarrows Y$ and $g:\subseteq Z\rightrightarrows W$
be multivalued functions on represented spaces.

\begin{itemize}
\item We say $f$ is \textit{Weihrauch reducible} to $g$ and write $f\leq_{\mathrm W} g$
if there exist computable partial functions $H,K:\subseteq \Baire \to \Baire$ such that
for every realizer $G\vdash g$, the function
\[
p \ \mapsto\ K\bigl(\pairWhite{p}{G(H(p))}\bigr)
\]
is a realizer of $f$.

\item We say $f$ is \textit{strongly Weihrauch reducible} to $g$ and write
$f\leq_{\mathrm{sW}} g$ if there exist computable $H,K$ such that for every $G\vdash g$,
\[
p \ \mapsto\ K\bigl(G(H(p))\bigr)
\]
realizes $f$ (so the post-processor $K$ does \textit{not} see the original input name $p$).
\end{itemize}

We write $f\equiv_{\mathrm W}g$ if $f\leq_{\mathrm W}g$ and $g\leq_{\mathrm W}f$.
\end{definition}

We will also use the notation of a \enquote{tightening}, which is defined as follows: For mulitvalued functions $f,g:\subseteq X \rightrightarrows Y$ on represented spaces $X,Y$, we call $f \preceq g$ a tightening of $g$ by $f$ iff $(\dom(g) \subseteq \dom(f) \, \land \, (\forall x \in \dom(g)) f(x) \subseteq g(x))$.

\begin{definition}[SCI-rank on the Weihrauch preorder]\label{def:rank}
For any (multi-valued) $f$ between represented spaces, define its \textit{Weihrauch-SCI rank}
\[
\mathrm{rank}_{\mathrm{W}}(f) \in \N_0\cup\{\infty\}
\]
by
\[
\mathrm{rank}_{\mathrm{W}}(f)
:=
\begin{cases}
\min\{k\in\N_0:\ f\le_{\mathrm{W}} \Limi{k}\}, & \text{if this set is nonempty,}\\
\infty, & \text{otherwise.}
\end{cases}
\]
\end{definition}

\begin{theorem}[Existence and well-posedness of $\mathrm{rank}_{\mathrm{W}}$]\label{thm:rank-wellposed}
For every (multi-valued) $f$ between represented spaces, $\mathrm{rank}_{\mathrm{W}}(f)$ is well-defined.
Moreover:
\begin{enumerate}
\item If $f\equiv_{\mathrm{W}} g$, then $\mathrm{rank}_{\mathrm{W}}(f)=\mathrm{rank}_{\mathrm{W}}(g)$ (so it descends to Weihrauch degrees).
\item If $f\le_{\mathrm{W}} g$, then $\mathrm{rank}_{\mathrm{W}}(f)\le \mathrm{rank}_{\mathrm{W}}(g)$ (order-preserving).
\end{enumerate}
\end{theorem}

\begin{proof}
\textbf{Well-definedness:} Let $S_f:=\{k\in\N_0:\ f\le_{\mathrm{W}}\Limi{k}\}$.
If $S_f=\emptyset$, define $\mathrm{rank}_{\mathrm{W}}(f):=\infty$.
If $S_f\neq\emptyset$, then $S_f\subseteq\N_0$ has a least element by well-ordering of $\N_0$; define $\mathrm{rank}_{\mathrm{W}}(f)$ to be this minimum.

\textbf{(1):} If $f\equiv_{\mathrm{W}} g$, then $f\le_{\mathrm{W}}g$ and $g\le_{\mathrm{W}}f$.
Hence $S_g\subseteq S_f$ and $S_f\subseteq S_g$, so $S_f=S_g$ and the minima (or emptiness) coincide.

\textbf{(2):} If $f\le_{\mathrm{W}}g$ and $k\in S_g$, then $g\le_{\mathrm{W}}\Limi{k}$.
By transitivity of $\le_{\mathrm{W}}$, $f\le_{\mathrm{W}}\Limi{k}$, so $k\in S_f$.
Thus $S_g\subseteq S_f$, hence $\min S_f \le \min S_g$ in the nonempty case, and the $\infty$ case is immediate.
\end{proof}

\begin{remark}[This is the \enquote{rank along the iterated-limit chain}]\label{rem:rank-chain}
The map $\mathrm{rank}_{\mathrm{W}}$ is a canonical order-preserving map from the Weihrauch preorder to the chain $(\N_0\cup\{\infty\},\le)$ induced by the increasing family
$\Limi{0}\le_{\mathrm{W}}\Limi{1}\le_{\mathrm{W}}\Limi{2}\le_{\mathrm{W}}\cdots$.
No further SCI-specific structure is needed for \cref{thm:rank-wellposed}.
\end{remark}

\subsection{Connecting \texorpdfstring{$\mathrm{rank}_{\mathrm{W}}$}{rankW} Back To SCI Towers (Type-2 SCI)}

\begin{definition}[Type-2 SCI for a represented-space problem]\label{def:type2sci}
Let $f:\subseteq X\rightrightarrows Y$ be a (multi-valued) function between represented spaces.
Define the \textit{Type-2 SCI} of $f$ by
\[
\mathrm{SCI}_{\mathrm{TTE}}(f):=\mathrm{rank}_{\mathrm{W}}(f).
\]
\end{definition}

\begin{remark}[Why this is the correct \enquote{effective} SCI notion]\label{rem:np-remark}
In the Type-2 setting, the fundamental discontinuity operator is $\Lim$.
The equality $\mathrm{SCI}_{\mathrm{TTE}}(f)\le k \iff f\le_{\mathrm{W}}\Limi{k}$
is exactly the content of \cref{def:type2sci}; it is commonly recorded as a theorem/observation in the literature, see e.g., \cite[Observation 43]{neumann2018topological} but the \textit{logical well-posedness} is already captured by \cref{thm:rank-wellposed}.
\end{remark}

\begin{definition}[Weihrauch-rank of a countable-$\Lambda$ SCI problem]\label{def:rank-of-SCIproblem}
Let $\mathcal{P}=\bigl( \Xi,\Omega,(\mathcal{M},d),\Lambda \bigr)$ be an SCI computational problem in the sense of \cref{def:comp-prob}-\cref{def:sci}
with \textit{countable} $\Lambda$ in the sense of \cref{def:countable-Lambda}.
Assume additionally that $\mathcal{M}$ is equipped with a fixed representation $\delta_{\mathcal{M}}$
(e.g.\ the Cauchy representation induced by $(\mathcal{M},d)$ if $\mathcal{M}$ is separable and complete).

Let $\widehat{\Xi}:\mathcal{I}_\Lambda\to\mathcal{M}$ be the factor map from \cref{lem:factor},
viewed as a (single-valued) function between represented spaces
$(\mathcal{I}_\Lambda,\delta_{\mathcal{I}_\Lambda})$ and $(\mathcal{M},\delta_{\mathcal{M}})$.

Define the \textit{Weihrauch-SCI rank} of the SCI problem $\mathcal{P}$ by
\[
\mathrm{rank}_{\mathrm{W}}(\mathcal{P})
:=\mathrm{rank}_{\mathrm{W}}(\widehat{\Xi}).
\]
\end{definition}

\begin{theorem}[Existence / well-posedness of $\mathrm{rank}_{\mathrm{W}}(\mathcal{P})$]\label{thm:rank-P}
Let $\mathcal{P}$ satisfy the hypotheses of \cref{def:rank-of-SCIproblem}.
Then $\mathrm{rank}_{\mathrm{W}}(\mathcal{P})$ exists and is independent of
\begin{enumerate}
\item the chosen enumeration of the countable set $\Lambda$, \textit{up to computable permutations
of $\N$} (i.e.\ up to computable reindexings of coordinates);
\item the choice of standard \textit{Cauchy} representation of $\C$
(with respect to the usual metric and a computable dense sequence);
\item the choice of representation of $\mathcal{M}$ \textit{within its computable isomorphism class}
(i.e.\ replacing $\delta_\mathcal{M}$ by a computably equivalent representation).
\end{enumerate}
\end{theorem}

\begin{proof}
Existence is immediate from \cref{thm:rank-wellposed} applied to the represented-space map $\widehat{\Xi}$.
Independence under computable reindexings follows from \cref{rem:enum-indep} and invariance
of $\mathrm{rank}_{\mathrm{W}}$ under computable isomorphisms.
Independence of standard admissible representations can be proven similarly: changing to a computably isomorphic representation
does not change the Weihrauch degree of $\widehat{\Xi}$, hence not its rank.
\end{proof}

\begin{remark}[Scope and limitations]\label{rem:scope}
\Cref{thm:rank-wellposed} and \cref{thm:rank-P} give a fully precise, foundational \textit{existence and well-posedness} theory for an SCI-rank compatible with Weihrauch theory. But one has to be careful: This is an \textit{effective} refinement of SCI.
The original SCI definitions (cf. \cref{def:comp-prob}-\cref{def:sci}) are not, by themselves,
statements about computability in the Type-2 sense; they must be enriched with representations and an admissible class of pre-/post-processing operations to connect to Weihrauch degrees.
\end{remark}

\begin{remark}[Classical status of $\mathrm{rank}_{\mathrm W}$ and $\mathrm{SCI}_{\mathrm{TTE}}$]
In Weihrauch theory, the limit operator $\lim$ is a central benchmark degree.
It characterizes limit computability: $f \le_{\mathrm W} \lim \iff f$ is limit computable.
Moreover, reducibility to iterates of $\lim$ corresponds to effective Borel measurability levels,
and $\lim$ enjoys strong algebraic properties (cylinder, strong parallelizability, etc.).
\end{remark}

\begin{remark}[Relation to SCI type $A$ and type $G$ in the SCI literature]
In the SCI framework, type $A$ towers impose effectivity restrictions at the base level,
whereas type $G$ towers are (typically) purely information-based and allow arbitrary
post-processing of finitely many oracle values without measurability/continuity constraints.

Therefore, $\mathrm{SCI}_G$ does \textit{not} automatically coincide with $\mathrm{SCI}_{\mathrm{TTE}}$.
For type $A$, one often expects compatibility with $\mathrm{SCI}_{\mathrm{TTE}}$ in analytic settings,
but this requires that the evaluation device provides Type-2 names/approximations of oracle values
and that the base computations avoid non-computable primitives (e.g.\ exact real comparisons).
\end{remark}

\begin{definition}[Cylinder (strong form)]
For a problem $g:\subseteq Z \rightrightarrows W$, let
\[
\mathrm{id}\times g:\subseteq X\times Z \rightrightarrows X\times W,
\qquad
(\mathrm{id}\times g)(x,z):=\{(x,w): w\in g(z)\}.
\]
We call $g$ a (strong) cylinder if
\[
\mathrm{id}\times g \leq_{sW} g.
\]
\end{definition}

\begin{definition}[Transparency (operational form on Baire)]\label{def:transparent}
A problem $g$ is \textit{transparent} if for every computable functional
$T:\subseteq\NN^\NN \to\NN^\NN$ there exists a computable functional
$S:\subseteq\NN^\NN \to\NN^\NN$ such that
\[
g \circ S \preceq T \circ g .
\]
\end{definition}

The facts that \(\Lim\) is a transparent cylinder and that iterated limits inherit this structure are standard in Weihrauch theory; see, for example,
\cite[Section~6, Theorem~6.5]{brattka2021weihrauch}.

\begin{lemma}[$\Limi{n}$ is a transparent cylinder]{\label{lem:limntranspcyl}}
For every $n\in\N_0$, the problem $\Limi{n}$ is a transparent cylinder.
\end{lemma}

\begin{proof}
This result is very classical, still we give a short proof here. We argue by induction on $n$.

\textbf{Base case $n=0$:} By definition, $\Limi{0}=\mathrm{id}_{\NN^{\NN}}$.
The identity is a cylinder since $\mathrm{id}\times \mathrm{id}\le_{\mathrm{sW}}\mathrm{id}$
via computable pairing/unpairing. It is transparent because for any computable $h$ we can take
$h':=h$ and obtain $h\circ \mathrm{id}=\mathrm{id}\circ h'$.

\textbf{Inductive step:} Assume $\Limi{n}$ is a transparent cylinder. Since $\Lim$ itself is a
transparent cylinder, and the class of transparent cylinders is closed under composition,
it follows that $\Lim\circ \Limi{n}=\Limi{n+1}$ is a transparent cylinder.

Thus $\Limi{n}$ is a transparent cylinder for all $n\in\NN_0$.
\end{proof}

\begin{lemma}[Normal form for reductions to a single-valued transparent cylinder]{\label{lem:TowerNormalForm}}
Let $g$ be a \textit{single-valued} transparent cylinder on Baire space, and let
\[
        f:\subseteq\mathbb N^{\mathbb N}\rightrightarrows\mathbb N^{\mathbb N}
\]
be a Baire-space problem. Then
\[
f \leq_{\mathrm W} g
\quad\Longleftrightarrow\quad
\exists\text{ computable }K:\subseteq\mathbb{N}^\mathbb{N}\to\mathbb{N}^\mathbb{N}
\text{ such that } g \circ K \preceq f.
\]
\end{lemma}

\begin{proof}
\textbf{($\Rightarrow$):}
Assume $f \leq_W g$. Since \(g\) is a cylinder, ordinary Weihrauch reducibility to \(g\) can be converted
to strong Weihrauch reducibility by the standard cylinder absorption argument; see \cite{brattka2021weihrauch}. Hence we may assume that the reduction is strong. Hence there exist
computable functionals $H_0,T$ such that
\[
T \circ g \circ H_0 \preceq f .
\]
By transparency of $g$, there exists a computable functional $S$ such that
\[
g \circ S \preceq T \circ g .
\]
Therefore
\[
g \circ S \circ H_0 \preceq T \circ g \circ H_0 \preceq f ,
\]
so we can take $K := S \circ H_0$.

\textbf{($\Leftarrow$):}
Suppose $g \circ K \preceq f$ for some computable $K$. Then
\[
\dom(f) \subseteq \dom(g \circ K)
\quad\text{and}\quad
(g \circ K)(x) \subseteq f(x)
\]
for all $x \in \dom(f)$. Hence the reduction is witnessed by the pre-processor $K$ and
the identity post-processor.
\end{proof}

We now state the corresponding name-level normal form. The next definition makes the typing precise: a represented-space problem is converted into an ordinary Baire-space
problem on names.

\begin{definition}[Name-level problem associated with a represented-space problem]\label{def:name-level-problem}
Let
\[
        f:\subseteq X\rightrightarrows Y
\]
be a problem between represented spaces
\[
        X=(|X|,\delta_X),\qquad Y=(|Y|,\delta_Y).
\]
The associated name-level problem is the multivalued map
\[
        \mathcal N_f:\subseteq\mathbb N^{\mathbb N}\rightrightarrows
        \mathbb N^{\mathbb N}
\]
defined by
\[
        \operatorname{dom}(\mathcal N_f) :=
        \{p\in\operatorname{dom}(\delta_X):
        \delta_X(p)\in\operatorname{dom}(f)\},
\]
and
\[
        \mathcal N_f(p) :=
        \{q\in\operatorname{dom}(\delta_Y):
        \delta_Y(q)\in f(\delta_X(p))\}.
\]
If \(f\) is single-valued, this means
\[
        \mathcal N_f(p)=
        \delta_Y^{-1}\bigl(\{f(\delta_X(p))\}\bigr).
\]
Whenever a formula of the form
\[
        \Limi{k} \circ K\preceq f
\]
appears with \(f\) a represented-space problem, it means
\[
        \Limi{k} \circ K\preceq \mathcal N_f
\]
as a statement between Baire-space problems.
\end{definition}

\begin{definition}[TTE tower of height $n$ for a problem]
Let $f:\subseteq X\rightrightarrows Y$ be a problem between represented spaces.
We say that $f$ has a \textit{TTE tower of height $n$} if
\[
f\ \le_{\mathrm W}\ \Limi{n}.
\]
\end{definition}

\begin{theorem}[Pure-tower normal form]\label{thm:pure-tower-normal-form}
Let $f$ be any problem between represented spaces. Then for every $n\in\N_0$
\[
f\le_{\mathrm W}\Limi{n}
\quad\Longleftrightarrow\quad
\exists\ \text{computable }K:\subseteq\Baire\to\Baire\ \text{ such that } \Limi{n} \circ K \preceq \mathcal N_f.
\]
In particular, the definition $\mathrm{SCI}_{\mathrm{TTE}}(f)=\mathrm{rank}_{\mathrm W}(f)$ coincides with the least height of a \enquote{computable $n$-limit normal form} for $f$.
\end{theorem}
\begin{proof}
\textbf{($\Rightarrow$):} By \cref{lem:limntranspcyl}, $\Limi{n}$ is a single-valued
transparent cylinder. Hence \cref{lem:TowerNormalForm} yields a computable $K$ such that
\[
\Limi{n} \circ K \preceq \mathcal N_f .
\]

\textbf{($\Leftarrow$):} If $\Limi{n} \circ K \preceq \mathcal N_f$ for some computable $K$, then
$f \leq_W \Limi{n}$ is witnessed by taking the Weihrauch pre-processor to be $K$
and the post-processor to be the identity.
\end{proof}

The preceding normal form is the first place where the distinction between raw SCI and
implemented SCI becomes visible. A raw tower may list its approximants separately; a
Weihrauch reduction must instead provide one name-level procedure producing the whole
input to the iterated limit operator. The next section shows that without such regularity
or uniformity assumptions, raw type-\(G\) SCI has very little Weihrauch content.

\begin{remark}[Uniformity of the whole approximation table]
The normal form in \cref{thm:pure-tower-normal-form} is stronger than merely
having, for each multi-index separately, a computable approximant. It requires a single
computable functional
\[
        K:\subseteq\mathbb N^{\mathbb N}\to\mathbb N^{\mathbb N}
\]
which, from one input name, produces the entire name to which the iterated limit
operator is applied. This is precisely the uniformity missing from non-uniform
deepest-level towers, as shown in \cref{prop:weak-Hansen-no-K}.
\end{remark}

\section{G-Case}{\label{sec:G-Case}}
\subsection{Type-G SCI (Unrestricted Post-Processing)}
We shortly recall the $G$-case definition of the SCI.

\begin{definition}[Raw type-G towers and \(\SCIG\)]\label{def:scig}
Let
\[
        \mathcal P=(\Xi,\Omega,(\mathcal M,d),\Lambda)
\]
be an SCI computational problem. A raw type-G tower is a standard tower in the sense
of \cref{def:standard-tower}; equivalently, only the deepest-level maps are
required to be general algorithms, and no computability, continuity, Borel, or uniform
implementation requirement is imposed on the indexed family. Define \(\SCIG(\mathcal P)\) as
in \cref{def:sci}.
\end{definition}

\begin{remark}[What type-G does not assume]\label{rem:typeg-not}
Raw type-\(G\) imposes no computability, continuity, Borel-measurability, or uniform implementation requirement on finite-transcript post-processing at the deepest level. This is not a defect in the definition; it is the point of the unrestricted information model. The consequences of this choice are made explicit in \cref{thm:finite-factor-scig0}.
\end{remark}

\subsection{Purely Logical Consequences Of The Type-G Setting}

\begin{definition}[Finite-query factorisation]\label{def:finite-factor}
Let $\mathcal{P}=(\Xi,\Omega,(\mathcal{M},d),\Lambda)$.
We say that $\Xi$ admits a \textit{finite-query factorisation through $\Lambda$} if there exist
$f_1,\dots,f_m\in\Lambda$ and a map $G: \operatorname{im}(f_1,\dots,f_m)\to \mathcal{M}$ such that
\[
\Xi(A)=G\bigl(f_1(A),\dots,f_m(A)\bigr)\qquad\forall A\in\Omega,
\]
where $\operatorname{im}(f_1,\dots,f_m)\subseteq \C^m$ is the range of $A\mapsto (f_1(A),\dots,f_m(A))$.
\end{definition}

\begin{theorem}[Finite-query factorisations force $\mathrm{SCI}_G=0$]\label{thm:finite-factor-scig0}
Let $\mathcal{P}=\bigl( \Xi,\Omega,(\mathcal{M},d),\Lambda \bigr)$ be an SCI computational problem. Suppose there exist
$f_1,\dots,f_m\in\Lambda$ and a function $G:\mathbb{C}^m\to \mathcal{M}$ such that
\[
\Xi(A)=G\bigl(f_1(A),\dots,f_m(A)\bigr)\quad\text{for all }A\in\Omega,
\]
i.e. $\Xi$ admits a finite-query factorisation through $\Lambda$.
Then $\mathrm{SCI}_G(\mathcal{P})=0$.
\end{theorem}
\begin{proof}
Define $\Gamma:\Omega\to \mathcal{M}$ by $\Gamma(A):=G(f_1(A),\dots,f_m(A))$ and define the query policy
$\Lambda_\Gamma(A):=\{f_1,\dots,f_m\}$ for all $A\in\Omega$.

Then $(\Gamma,\Lambda_\Gamma)$ is a general algorithm: if $B$ agrees with $A$ on all queries in
$\Lambda_\Gamma(A)$, then $f_i(B)=f_i(A)$ for all $i$, hence $\Gamma(B)=\Gamma(A)$, and clearly also
$\Lambda_\Gamma(B)=\Lambda_\Gamma(A)$ because the query set is constant.

Moreover $\Gamma=\Xi$ by construction, so $\Xi$ is computed by a height-$0$ tower (i.e.\ a single
general algorithm). Therefore $\mathrm{SCI}_G(\mathcal{P})=0$.
\end{proof}

\begin{remark}[Core logical point]\label{rem:core-logical}
\Cref{thm:finite-factor-scig0} is \textit{pure logic}: once base post-processing is unrestricted,
\textit{any} map $G$ on the finite transcript is allowed.
This is the reason why descriptive-set/Type-2 lower-bound methods do not automatically imply type-G lower bounds.
\end{remark}

\subsection{Monotonicity In Tower Type: G Vs. A}

\begin{theorem}[Monotonicity in tower type]\label{thm:monotone-G-vs-BSS-A}
For every SCI computational problem $\mathcal P$ for which $\SCIAbss(\mathcal P)$
is defined, we have
\[
\mathrm{SCI}_{G}(\mathcal P)\ \le\ \SCIAbss(\mathcal P).
\]
\end{theorem}
\begin{proof}
Let $k:=\SCIAbss(\mathcal P)$ and fix a height-$k$ arithmetic tower
witnessing this value. Each base map of this tower is a BSS-arithmetic algorithm
in the sense of \cref{def:BSS-arith-alg}, hence, in particular, it depends on only finitely many oracle
values and is therefore a general algorithm in the sense of \cref{def:gen-alg}.
Thus the same tower is admissible as a type-$G$ tower, witnessing
$\mathrm{SCI}_{G}(\mathcal P)\le k$.
\end{proof}

\subsection{Why Type-G Is Not Captured By Weihrauch/TTE In General}

\begin{lemma}[Finite $\mathrm{rank}_{\mathrm W}$ implies Borel regularity]\label{lem:borel-regularity}
Let $f:\subseteq X\to Y$ be single-valued between represented spaces and assume
$f\le_{\mathrm W}\Limi{n}$ for some $n \in \NN_0$.
Then $f$ is Borel-measurable with respect to the represented-space induced topologies.
In particular, if $Y=\{0,1\}$ is discrete, then $f^{-1}(\{1\})$ is a Borel subset of $X$.
\end{lemma}

\begin{proof}
If $n=0$, then $\Limi{0}=\mathrm{id}$ and $f\le_{\mathrm W}\mathrm{id}$ implies that $f$ is computable,
hence continuous, hence Borel measurable.

Now assume $n\in\NN$. By the characterization theorem for the $\Lim$-hierarchy (c.f.\ \cite[Thm.~6.5]{brattka2021weihrauch}),
$f\le_{\mathrm W}\Limi{n}$ implies that $f$ is effectively $\Sigma^0_{n+1}$-measurable, and in particular
Borel measurable.
\end{proof}

\begin{theorem}[Type-G can have $\mathrm{SCI}_G=0$ while $\mathrm{SCI}_{\mathrm{TTE}}=\infty$]\label{thm:scig0-but-tteinf}
Assume a standard represented space structure on $X=\C$ (e.g.\ Cauchy representation),
and let $Y=\{0,1\}$ be discrete.
Let $A\subseteq \C$ be a fixed analytic-complete (hence non-Borel) set.
Define the single-valued map
\[
f:\C\to\{0,1\},\qquad f(z):=\mathbf{1}_A(z).
\]
Consider the SCI computational problem
\[
\mathcal{P}=(\Xi,\Omega,(\mathcal{M},d),\Lambda)
\quad\text{with}\quad
\Omega=\C,\ \mathcal{M}=\{0,1\},\ \Xi=f,\ \Lambda=\{\mathrm{id}_\C\}.
\]
Then
\begin{enumerate}
\item $\mathrm{SCI}_G(\mathcal{P})=0$;
\item $\mathrm{SCI}_{\mathrm{TTE}}(f)=\infty$ (equivalently, $\mathrm{rank}_{\mathrm W}(f)=\infty$).
\end{enumerate}
\end{theorem}

\begin{proof}
We note first of all that such an analytic-complete set $A$ indeed exists, see, for example \cite[Sec.~25]{Kec95}.\\
\textbf{(1):} Since $f=\mathbf{1}_A\circ \mathrm{id}_\C$, this is a finite-query factorisation
through the single evaluation $\mathrm{id}_\C\in\Lambda$ with post-processing $G=\mathbf{1}_A$.
Thus $\mathrm{SCI}_G(\mathcal{P})=0$ by \cref{thm:finite-factor-scig0}.

\textbf{(2):} Suppose for contradiction that $f\le_{\mathrm W}\Limi{n}$ for some $n$.
Then by \cref{lem:borel-regularity}, the set $A=f^{-1}(\{1\})$ must be Borel,
contradicting the choice of $A$. Hence $f\not\le_{\mathrm W}\Limi{n}$ for all $n$, so $\mathrm{SCI}_{\mathrm{TTE}}(f)=\infty$.
\end{proof}

\begin{corollary}[No bound from raw type-G height to Weihrauch rank]\label{cor:no-bound-G-to-W}
There is no function
\[
        h:\mathbb N_0\cup\{\infty\}\to\mathbb N_0\cup\{\infty\}
\]
such that, for every countable-\(\Lambda\) SCI computational problem equipped with an arbitrary represented-space enrichment,
\[
        \SCIG(\mathcal P)\le n \quad\Longrightarrow\quad \mathrm{rank}_{\mathrm W} (\widehat{\Xi})\le h(n).
\]
\end{corollary}

\begin{proof}
\Cref{thm:scig0-but-tteinf} gives a problem with $\SCIG(\mathcal P)=0$ but
\[
        \mathrm{rank}_{\mathrm W} (\widehat{\Xi})=\infty.
\]
Thus no finite value of \(h(0)\) can make the implication true.
\end{proof}

\begin{remark}[The reverse direction requires a compatibility package]
A converse implication from Weihrauch rank to raw type-G SCI height is not a well-typed general theorem unless one fixes a compatibility package relating the
represented input names to the finite-evaluation interface \(\Lambda\). In compatible effective settings such implications must be proved from additional assumptions. This
is exactly the role of implemented towers and of the pure $\mathcal R-\Lim$ package in \cref{sec:uniformity-minimality}.
\end{remark}

The preceding results show that raw type-\(G\) SCI is too permissive for
descriptive-set-theoretic lower bounds. To recover such lower bounds one must impose
regularity at the deepest level of the tower. This motivates the regularity-restricted
hierarchy developed in Section~\ref{sec:IntHierarchy}.

\section{The Intermediate Hierarchy}{\label{sec:IntHierarchy}}
The last section is the inspiration point for our full intermediate hierarchy now. But first of all we sharply precise what we mean by a hierarchy.
\begin{definition}[Preorder]
A \textit{preorder} is a pair $(I,\preceq)$ where $I$ is a set and $\preceq$ is a relation on $I$
that is reflexive and transitive.
\end{definition}

\begin{definition}[Hierarchy indexed by a preorder]\label{def:hierarchy}
Let $U$ be a set (the \enquote{universe}) and let $(I,\preceq)$ be a preorder.
A family of subclasses $(\mathcal{H}_i)_{i\in I}$ with $\mathcal{H}_i\subseteq U$ is called an
\textit{$(I,\preceq)$-hierarchy on $U$} if
\[
i\preceq j \quad\Longrightarrow\quad \mathcal{H}_i \subseteq \mathcal{H}_j.
\]
\end{definition}

\begin{definition}[Regularity classes on $\Baire$]\label{def:reg-classes}
Let $\mathcal{R}$ be a class of partial maps $F:\subseteq \Baire\to\Baire$.
We consider the following standard classes (all w.r.t.\ the usual product topology on $\Baire$):
\begin{itemize}
\item $\Comp$: the (Type-2) computable partial maps.
\item $\Cont$: the continuous partial maps.
\item $\Bor$: the Borel-measurable partial maps (measurable w.r.t.\ the Borel $\sigma$-algebra on $\Baire$).
\item $\BaireMeas$: the Baire-measurable partial maps (measurable w.r.t.\ the $\sigma$-algebra generated by the
zero sets / equivalently by continuous real-valued maps).
\end{itemize}
We order such classes by inclusion: $\mathcal{R}\subseteq\mathcal{S}$.
\end{definition}

\begin{proposition}[Inclusion chain]\label{prop:reg-inclusions}
One has
\[
\Comp \subseteq \Cont \subseteq \Bor \subseteq \BaireMeas.
\]
\end{proposition}

\begin{proof}
On Baire space \(\mathbb N^{\mathbb N}\), every Type-2 computable partial functional is continuous on its domain. Every continuous partial map is Borel measurable, and every Borel measurable partial map is Baire measurable because every Borel set in a Polish space has the Baire property. This proves the inclusion chain.
\end{proof}

\subsection{Continuous And Borel-Weihrauch Reducibilities}{\label{sec:ContBorelWeihrauchRed}}
\subsubsection{Parametrized Weihrauch Reducibilities}

\begin{definition}[$\mathcal{R}$-Weihrauch reducibility]\label{def:R-weihrauch}
Let $f:\subseteq X\rightrightarrows Y$ and $g:\subseteq U\rightrightarrows V$ be problems between represented spaces.
Let $\mathcal{R}$ be a class of partial maps $\subseteq\Baire\to\Baire$.
We write
\[
f \le_{\mathrm W}^{\mathcal{R}} g
\]
if there exist $H,K\in \mathcal{R}$ such that for every realizer $G\vdash g$,
the map
\[
p \mapsto K\bigl(\langle p, G(H(p))\rangle\bigr)
\]
is a realizer of $f$. (Here $\langle\cdot,\cdot\rangle:\Baire\times\Baire\to\Baire$ is any fixed computable pairing.)
\end{definition}

\begin{remark}[Standard special cases]
\hspace{1mm}
\begin{itemize}
\item $\le_{\mathrm W}^{\Comp}$ is the usual (computable) Weihrauch reducibility.
\item $\le_{\mathrm W}^{\Cont}$ is the standard \textit{continuous Weihrauch reducibility}: Weihrauch reducibility has a continuous counterpart obtained by replacing the reduction functions $H,K$ by continuous ones, see e.g. \cite{brattka2011weihrauch}.
\end{itemize}
\end{remark}

\begin{proposition}[$\le_{\mathrm W}^{\mathcal{R}}$ is a preorder and monotone in $\mathcal{R}$]\label{prop:R-weihrauch-preorder}
Fix $\mathcal{R}$ that is closed under composition and contains all computable pairing/projection maps on $\Baire$.
Then $\le_{\mathrm W}^{\mathcal{R}}$ is a preorder.
Moreover, if $\mathcal{R}\subseteq\mathcal{S}$, then
\[
f \le_{\mathrm W}^{\mathcal{R}} g \ \Longrightarrow\ f \le_{\mathrm W}^{\mathcal{S}} g.
\]
\end{proposition}

\begin{proof}
Reflexivity is witnessed by taking $H(p)=p$ and $K(\langle p,q\rangle)=q$, which are computable (hence in $\mathcal{R}$).
Transitivity is obtained by composing witnesses; closure of $\mathcal{R}$ under composition and pairing ensures
the composed witnesses remain in $\mathcal{R}$.
Monotonicity in $\mathcal{R}$ is immediate: any witnesses in $\mathcal{R}$ are also witnesses in $\mathcal{S}$.
\end{proof}

\begin{definition}[Continuous and Borel Weihrauch reducibilities]\label{def:cont-borel-weihrauch}
Define
\[
\le_{\mathrm W}^{\mathrm{cont}} := \le_{\mathrm W}^{\Cont},
\qquad
\le_{\mathrm W}^{\mathrm{Bor}} := \le_{\mathrm W}^{\Bor}.
\]
\end{definition}

\begin{remark}[Oracle-relative view of $\le_{\mathrm W}^{\mathrm{cont}}$]
A common equivalent presentation is: $f\le_{\mathrm W}^{\mathrm{cont}} g$ iff there exists an oracle
$p\in\Baire$ such that $f\le_{\mathrm W}^{p} g$ (i.e.\ the reduction functionals are computable relative to $p$).
This is exactly the viewpoint used when Brattka-Gherardi-Pauly relativize \cite[Theorem 6.5]{brattka2021weihrauch} and characterize the non-effective Borel hierarchy via existence of an oracle.
\end{remark}

\subsection{Restricted SCI Towers}{\label{sec:RestrSCITowers}}

\subsubsection{Topological Enrichment}

\begin{definition}[Topologized SCI problem]\label{def:topoSCI}
A \textit{topologized SCI computational problem} is a tuple
\[
\mathcal{P}_\tau=(\Xi,\Omega,(\mathcal{M},d),\Lambda,\tau)
\]
where $(\Xi,\Omega,(\mathcal{M},d),\Lambda)$ is a computational problem in the SCI sense, and $\tau$ is a topology on $\Omega$.
We write $\mathcal{B}(\Omega)$ for the Borel $\sigma$-algebra of $(\Omega,\tau)$ and $\mathcal{B}(\mathcal{M})$ for the Borel $\sigma$-algebra of the metric space $(\mathcal{M},d)$.
\end{definition}

\begin{remark}[Why the extra structure is necessary]
To state that a subset \(A\subseteq\Omega\) is Borel, analytic, or complete for a
descriptive pointclass, one must equip \(\Omega\) with a topology or at least a standard
Borel structure. Thus any attempt to import descriptive-set-theoretic lower bounds
into SCI necessarily fixes such additional structure.
\end{remark}

\subsubsection{Regularity-Restricted Base Algorithms And Towers}
Let
\[
\mathfrak{R}_{\mathrm{top}} := \{\Cont,\Bor,\BaireMeas\}.
\]
For a topological space $X$ and a metric space $Y$, define
\[
\Reg_{\Cont}(X,Y) := \{h:X\to Y : h \text{ is continuous}\},
\]
\[
\Reg_{\Bor}(X,Y) := \{h:X\to Y : h \text{ is Borel measurable}\},
\]
\[
\Reg_{\BaireMeas}(X,Y) := \{h:X\to Y : h \text{ is Baire measurable}\}.
\]
We order the labels in $\mathfrak{R}_{\mathrm{top}}$ by
\[
r \sqsubseteq_{\mathrm{top}} s
\quad:\Longleftrightarrow\quad
\Reg_r(X,Y)\subseteq \Reg_s(X,Y)
\]
for every topological space $X$ and every metric space $Y$.
Equivalently,
\[
\Cont\sqsubseteq_{\mathrm{top}}\Bor
\sqsubseteq_{\mathrm{top}}\BaireMeas.
\]
Therefore in \cref{sec:RestrSCITowers} - \cref{sec:ExampleSCI} we use $r,s \in \mathfrak{R}_{\mathrm{top}}$ for topological regularity labels, other than we did in \cref{sec:ContBorelWeihrauchRed}, where we reserved calligraphic letters $\mathcal R,\mathcal S$ for classes of \textit{realizers} on $\mathbb{N}^{\mathbb{N}}$.

\begin{definition}[Fixed-query $(\Lambda,r)$-base algorithms]\label{def:fixed-query-base}
Let $\mathcal{P}_\tau=(\Xi,\Omega,(\mathcal{M},d),\Lambda,\tau)$ be a topologized SCI computational problem and let
$r \in \mathfrak{R}_{\mathrm{top}}$.
A general algorithm $\Gamma:\Omega\to \mathcal{M}$ is called \textit{fixed-query $(\Lambda,r)$-regular} if there exist $f_1,\dots,f_m\in\Lambda$ and a map $G:\C^m\to \mathcal{M}$ with $G\in \Reg_r(\mathbb{C}^m,\mathcal M)$ such that
\[
\Gamma(A)=G\bigl(f_1(A),\dots,f_m(A)\bigr)\quad\forall A\in\Omega,
\]
and each $f_i:(\Omega,\tau)\to\C$ belongs to $\Reg_r(\Omega,\mathbb{C})$.
\end{definition}

\begin{definition}[$r$-general algorithms and $r$-towers]\label{def:R-tower}
Let $\mathcal{P}_\tau=(\Xi,\Omega,(\mathcal{M},d),\Lambda,\tau)$ be a topologized SCI computational problem,
and let $r \in \mathfrak{R}_{\mathrm{top}}$.

An \textit{$r$-general algorithm} is a general algorithm $\Gamma:\Omega\to\mathcal{M}$ such that 
\[
\Gamma \in \Reg_r(\Omega,\mathcal M).
\]
Here, unlike \cref{def:fixed-query-base}, the finite query set may depend on the input!

A \textit{height-$k$ $r$-tower} is a standard tower in the sense of \cref{def:standard-tower} whose deepest-level maps $\Gamma_{n_k,\dots,n_1}:\Omega\to\mathcal{M}$ are $r$-general algorithms.
\end{definition}

The intuition here is: there are two distinct ways to enforce regularity at the base level.
\Cref{def:fixed-query-base} restricts the \textit{post-processing} from finitely many queried oracle values
(fixed-query form), while \cref{def:R-tower} restricts the \textit{entire} base algorithm as a map
$(\Omega,\tau)\to (\mathcal{M},d)$ (allowing adaptive querying but controlling descriptive complexity).
The restricted heights $\mathrm{SCI}_{G,r}$ quantify how many iterated limits remain
once this regularity is enforced.

\begin{definition}[Restricted SCI values]\label{def:SCI-R}
For $r\in\mathfrak{R}_{\mathrm{top}}$ define
\[
\mathrm{SCI}_{G,r}(\mathcal{P}_\tau)
\;:=\;
\min\Bigl\{
k\in\mathbb N_0 :
\text{there exists a height-}k\text{ $r$-tower computing }\Xi
\Bigr\},
\]
with value $\infty$ if the set is empty.
\end{definition}

\subsubsection{The Hierarchy Property}

\begin{theorem}[Restricted SCI models form a hierarchy in the regularity parameter]\label{thm:hierarchy-in-R}
Let $\mathcal{P}_\tau$ be a fixed topologized SCI problem.
If $r \sqsubseteq_{\mathrm{top}} s$,
then
\[
\mathrm{SCI}_{G,s}(\mathcal{P}_\tau) \ \le\ \mathrm{SCI}_{G,r}(\mathcal{P}_\tau).
\]
In particular,
\[
\mathrm{SCI}_{G,\BaireMeas}(\mathcal{P}_\tau)\ \le\ \mathrm{SCI}_{G,\Bor}(\mathcal{P}_\tau)\ \le\ \mathrm{SCI}_{G,\Cont}(\mathcal{P}_\tau).
\]
\end{theorem}
\begin{proof}
Any height-$k$ $r$-tower is in particular a height-$k$ $s$-tower whenever
$r \sqsubseteq_{\mathrm{top}} s$. The displayed chain follows from
\[
\Cont\sqsubseteq_{\mathrm{top}}\Bor
\sqsubseteq_{\mathrm{top}}\BaireMeas.
\]
\end{proof}

\begin{theorem}[A two-parameter hierarchy: regularity and height]\label{thm:two-parameter-hierarchy}
Fix a universe $U$ of topologized SCI computational problems.
For each pair $(r,k)$ with $r\in\mathfrak{R}_{\mathrm{top}}$ and $k\in\N_0$, define
\[
\mathcal{H}_{(r,k)} := \{\ \mathcal{P}_\tau\in U\ :\ \mathrm{SCI}_{G,r}(\mathcal{P}_\tau)\le k\ \}.
\]
Equip the index set
\[
I:=\mathfrak{R}_{\mathrm{top}}\times\mathbb N_0
\]
with the preorder
\[
(r,k)\preceq(s,\ell)
\quad:\Longleftrightarrow\quad
r\sqsubseteq_{\mathrm{top}} s \text{ and } k\le \ell.
\]
Then $(\mathcal{H}_{(r,k)})_{(r,k)\in I}$ is an $(I,\preceq)$-hierarchy in the sense of \cref{def:hierarchy}.
\end{theorem}

\begin{proof}
If $(r,k)\preceq(s,\ell)$, then any $r$-tower of height $\le k$
is also an $s$-tower and has height $\le \ell$. Hence $\mathcal{H}_{(r,k)}\subseteq \mathcal{H}_{(s,\ell)}$.
\end{proof}

\begin{definition}[Optional refinement: add a query-policy parameter]\label{def:Q-R-refinement}
Let $\mathcal{Q}:=\{\mathrm{fix},\mathrm{ad}\}$.
For $q=\mathrm{fix}$ use fixed-query $(\Lambda,r)$-regular base maps in the sense of \cref{def:fixed-query-base};
for $q=\mathrm{ad}$ use $r$-general base maps in the sense of \cref{def:R-tower}.

Define the corresponding base-map class
\[
\mathsf{Base}(q,r)
:=
\begin{cases}
\{\text{fixed-query $(\Lambda,r)$-regular base maps}\}, & q=\mathrm{fix},\\
\{\text{$r$-general base maps}\}, & q=\mathrm{ad}.
\end{cases}
\]

Define \(\mathrm{SCI}^{q,r}_G(\mathcal P_\tau)\) as the least height of a standard tower for \(\mathcal P_\tau\) whose deepest-level maps lie in \(\mathsf{Base}(q,r)\).
\end{definition}

\begin{definition}[Index preorder for the $(q,r)$-refinement]\label{def:QRK-preorder}
Let
\[
I := \mathcal Q \times \mathfrak{R}_{\mathrm{top}} \times \mathbb N_0.
\]
For $i=(q,r,k)$ and $j=(q',r',k')$ in $I$, define
\[
i\preceq j
\;:\Longleftrightarrow\;
\mathsf{Base}(q,r)\subseteq \mathsf{Base}(q',r')
\ \text{ and }\ k\le k'.
\]
Then $(I,\preceq)$ is a preorder.
\end{definition}

\begin{proposition}[Hierarchy property]\label{prop:QRK-hierarchy}
Fix a universe $U$ of topologized SCI problems.
For each index $i=(q,r,k)\in I$ define
\[
\mathcal{H}_i
\;:=\;
\bigl\{\mathcal{P}_\tau\in U : \mathrm{SCI}^{q,r}_G(\mathcal{P}_\tau)\le k\bigr\}.
\]
Then the family $(\mathcal{H}_i)_{i\in I}$ is an $(I,\preceq)$-hierarchy on $U$ in the sense of
\cref{def:hierarchy}.
\end{proposition}

\begin{proof}
Let $i=(q,r,k)\preceq j=(q',r',k')$ and let $\mathcal{P}_\tau\in\mathcal{H}_i$.
Then $\mathrm{SCI}^{q,r}_G(\mathcal{P}_\tau)\le k$, so there exists a tower of height $\le k$
computing $\mathcal{P}_\tau$ whose base maps lie in $\mathsf{Base}(q,r)$.
Since $\mathsf{Base}(q,r)\subseteq \mathsf{Base}(q',r')$ and $k\le k'$,
the same tower is admissible for $(q',r')$ and has height $\le k'$.
Hence $\mathrm{SCI}^{q',r'}_G(\mathcal{P}_\tau)\le k'$, i.e.\ $\mathcal{P}_\tau\in\mathcal{H}_j$.
Therefore $\mathcal{H}_i\subseteq \mathcal{H}_j$.
\end{proof}

\begin{remark}[Relation to Kihara-style reducibilities]
Kihara's piecewise-continuity reducibilities in \cite{kihara2017borel} provide a related regularity-sensitive
framework for reducibility and uniformization problems. The present hierarchy differs
in that every SCI problem comes with an explicit finite-information interface
\(\Lambda\), a tower-height parameter, and a distinction between fixed-query and
adaptive-query base policies. A more detailed comparison is given in the extended
arXiv version.
\end{remark}

\subsection{Comparison Theorems: What Each Restriction Buys You}

\subsubsection{Borel-Restricted Towers}

We shortly recall a standard but here very crucial result.

\begin{lemma}[Pointwise limits of Borel maps are Borel]\label{lem:pointwise-limit-borel}
Let \((\Omega,\mathcal A)\) be a measurable space and let \((\mathcal M,d)\) be a metric space.
Let
\[
        f_n:\Omega\to \mathcal M
\]
be Borel measurable maps, and suppose
\[
        f(\omega):=\lim_{n\to\infty}f_n(\omega)
\]
exists for every \(\omega\in\Omega\). Then \(f:\Omega\to \mathcal M\) is Borel measurable.
\end{lemma}

\begin{proof}
Let \(U\subseteq \mathcal M\) be open. For \(r>0\), define
\[
        U_r:=\{y\in \mathcal M:\operatorname{dist}(y, \mathcal M\setminus U)>r\},
\]
with the convention \(\operatorname{dist}(y,\varnothing)=+\infty\). Then
\[
        f^{-1}(U) =
        \bigcup_{m=1}^{\infty}
        \bigcup_{N=1}^{\infty}
        \bigcap_{n\ge N} f_n^{-1}(U_{2^{-m}}).
\]
The right-hand side is measurable, and the equality follows from the definition
of pointwise convergence and the openness of \(U\). Hence \(f\) is Borel measurable.
\end{proof}

\begin{theorem}[Borel towers compute only Borel targets]\label{thm:borelTowersBorelTargets}
Let
\[
\mathcal P_\tau=(\Xi,\Omega,(\mathcal M,d),\Lambda,\tau)
\]
be a topologized SCI problem with \(\tau\) Polish. If \(\Xi\) is computed by a
height-\(k\) tower whose deepest-level maps are Borel measurable as maps
\[
(\Omega,\tau)\to (\mathcal M,d),
\]
then
\[
\Xi:(\Omega,\tau)\to(\mathcal M,d)
\]
is Borel measurable.  In particular, for a decision problem \(\mathcal M=\{0,1\}\) with the
discrete topology, the acceptance set is Borel.
\end{theorem}

\begin{proof}
We prove the more general statement that the target map of every finite-height Borel tower into a metric space is Borel.

For height \(0\), the target map is itself a Borel deepest-level map.

Assume the result holds for heights \(<k\). If \(\Xi\) is computed by a height-\(k\)
Borel tower, then
\[
\Xi(\omega)=\lim_{n_k\to\infty} \Gamma_{n_k}(\omega)
\]
for maps \(\Gamma_{n_k}\) each computed by a height-\((k-1)\) Borel tower. By the
induction hypothesis, every \(\Gamma_{n_k}\) is Borel. By \cref{lem:pointwise-limit-borel}, the pointwise limit \(\Xi\) is Borel.
The decision-case statement follows by taking the preimage of \(\{1\}\).
\end{proof}

\begin{remark}[Why this recovers standard lower-bound methods]
\Cref{thm:borelTowersBorelTargets} is the formal premise behind many
descriptive-set-theoretic lower bounds: if all deepest-level maps are effective, recursive,
Type-2 computable, continuous, or Borel, then they are in particular Borel in the
relevant represented/topological structure. Hence a non-Borel acceptance set rules
out every finite-height tower of that restricted kind. The conclusion does not apply
to unrestricted type-G towers, because raw type-G imposes no Borel requirement on deepest-level finite-transcript post-processing.
\end{remark}

\subsubsection{Continuous Base: Baire-Class Height}

\begin{definition}[Baire classes (finite levels)]\label{def:baire-classes}
Let $X$ be a topological space and $Y$ a metric space.
Define inductively
\begin{itemize}
\item $\mathsf{BC}_0(X,Y)$ is the class of continuous maps $X\to Y$.
\item $\mathsf{BC}_{n+1}(X,Y)$ is the class of pointwise limits of sequences of maps in $\mathsf{BC}_n(X,Y)$.
\end{itemize}
\end{definition}

\begin{theorem}[Continuous-base towers yield finite Baire class]\label{thm:cont-tower-baireclass}
Let $\mathcal{P}_\tau$ be a topologized SCI problem.
If $\Xi$ is computed by a height-$k$ $\Cont$-tower in the sense of \cref{def:R-tower} with $r=\Cont$,
then $\Xi\in \mathsf{BC}_k(\Omega,\mathcal{M})$. If moreover $(\Omega,\tau)$ is Polish, then $\Xi$ is Borel measurable.
\end{theorem}

\begin{proof}
We argue by induction on $k$.

For $k=0$, $\Xi$ is a continuous general algorithm, hence $\Xi\in\mathsf{BC}_0$ by \cref{def:baire-classes}.

Assume the claim for $k-1$ and let $\Xi$ be computed by a height-$k$ $\Cont$-tower.
Unwinding the tower definition, there exist maps $\Gamma_{n_k}:\Omega\to\mathcal{M}$ with
\[
\Xi(\omega)=\lim_{n_k\to\infty}\Gamma_{n_k}(\omega)
\quad\text{for all }\omega\in\Omega,
\]
and each $\Gamma_{n_k}$ is itself computed by a height-$(k-1)$ $\Cont$-tower (the inner limits).
By the induction hypothesis, each $\Gamma_{n_k}\in \mathsf{BC}_{k-1}(\Omega,\mathcal{M})$.
Hence $\Xi$ is a pointwise limit of a sequence in $\mathsf{BC}_{k-1}$, i.e.\ $\Xi\in \mathsf{BC}_k$.

If moreover $(\Omega,\tau)$ is Polish, then every finite Baire class map is Borel measurable,
hence $\Xi$ is Borel measurable.
\end{proof}

The preceding results explain the logical effect of imposing regularity at the deepest
level of an SCI tower. Borel regularity recovers descriptive-set lower-bound methods,
while continuous regularity locates the target map in a finite Baire class. The next
section shows that the raw hierarchy itself is nontrivial: there are canonical source
problems of every finite raw type-\(G\) height.

\subsubsection{Continuous Weihrauch To \texorpdfstring{$\Limi{n}$}{limn} And Borel Levels}

\begin{theorem}[Continuous reducibility to $\Limi{n}$ characterizes Borel measurability level]\label{thm:bgp-relativized}
Let \(f\) be a single-valued problem between represented spaces for which the usual
notions of effective and relativized \(\Sigma^0_{n+1}\)-measurability are defined, for
instance between represented Polish spaces or admissibly represented countably based
spaces. Let \(n\in\mathbb N_0\).
\begin{enumerate}
\item $f\le_{\mathrm W}\Limi{n}$ if and only if $f$ is \textit{effectively} $\Sigma^0_{n+1}$-measurable.
\item For an oracle $p\in\mathbb{N}^{\mathbb{N}}$, $f\le_{\mathrm W}^{\,p}\Limi{n}$ if and only if
$f$ is $\Sigma^0_{n+1}$-measurable \textit{relative to $p$}.
\item Equivalently, $f\le_{\mathrm W}^{\mathrm{cont}}\Limi{n}$ if and only if $f$ is (non-effective)
$\Sigma^0_{n+1}$-measurable.
\end{enumerate}

At level $n=0$ one has
\[
f \le_{\mathrm W} \Limi{0}
\quad\Longleftrightarrow\quad
f \le_{\mathrm W} \mathrm{id}
\quad\Longleftrightarrow\quad
f \text{ is computable},
\]
and
\[
f \le_{\mathrm W}^{\mathrm{cont}} \Limi{0}
\quad\Longleftrightarrow\quad
f \le_{\mathrm W}^{\mathrm{cont}} \mathrm{id}
\quad\Longleftrightarrow\quad
f \text{ is continuous}.
\]
\end{theorem}
\begin{proof}
The case $n\!=\!0$ is immediate from $\Limi{0}=\mathrm{id}$. (1)-(2) is exactly \cite[Theorem~6.5]{brattka2021weihrauch}, together with its standard oracle relativization.
(3) follows from the well-known equivalence between continuous Weihrauch reducibility and
computable reducibility relative to some oracle.
\end{proof}

In summary we can now say

\begin{itemize}
\item The \textit{continuous} side is robust: \enquote{height $n$ with continuous base} forces (at most) Baire class $n$ (\cref{thm:cont-tower-baireclass}), and in the standard single-valued represented-Polish setting, continuous Weihrauch reducibility to \(\Limi{n}\) characterizes exactly the corresponding $\Sigma^0_{n+1}$-measurability level (\cref{thm:bgp-relativized}).

\item The \textit{Borel} restriction captures the minimal premise behind many lower bounds:
Borel base maps imply Borel targets (\cref{thm:borelTowersBorelTargets}), hence non-Borel acceptance sets block \textit{all} finite-height Borel/effective towers.

\item In contrast, the unrestricted type-G model drops the Borel/continuous premise.
Therefore descriptive-set hardness alone does not imply any type-G lower bound; see \cref{thm:finite-factor-scig0} and \cref{thm:scig0-but-tteinf}.
\end{itemize}

Thus the \enquote{intermediate hierarchy} is mathematically concrete: it is the two-parameter hierarchy in \cref{thm:two-parameter-hierarchy}, and its interaction with Weihrauch theory is governed by the measurability characterizations around $\Limi{n}$ (cf. \cite[Theorem 6.5]{brattka2021weihrauch}). For implemented SCI readings, the logical role of this hierarchy is analyzed in \cref{sec:uniformity-minimality}: raw type-\(G\) is too permissive, all-level finite-information towers are too restrictive, and the represented pure \(\mathcal R\)-\(\Lim\) normal form supplies the correct Weihrauch-comparable uniformity condition.

\section{Canonical Source Problems For Finite Type-G Heights}{\label{sec:ExampleSCI}}

This section gives a source family realizing arbitrary finite heights for the standard
raw type-G SCI of \cref{def:standard-tower}. These are raw extensional height statements, not claims of Weihrauch or Type-2 computability. The source problems serve two purposes. First, they show that the standard raw type-\(G\) hierarchy has nontrivial finite levels. Second, they provide
canonical hard sources for later finite-query transport arguments in concrete analytic or operator-theoretic settings.

\subsection{The Cantor-Matrix Input Space}

\begin{definition}[Cantor-matrix source space]
Let $\Omega_1:=2^{\mathbb N\times\mathbb N}$ equipped with the product topology. For \((a,b)\in\mathbb N^2\), define the coordinate evaluation
\[
\operatorname{ev}_{a,b}:\Omega_1\to\{0,1\}, \,
\operatorname{ev}_{a,b}(A):=A(a,b).
\]
Let
\[
\Lambda_{\mathrm{mat}} := \{\operatorname{ev}_{a,b}:(a,b)\in\mathbb N^2\}.
\]
For each \(m\in \mathbb{N}\), fix a bijection
\[
\beta_m:\mathbb N^m\to\mathbb N^2.
\]
\end{definition}

\begin{definition}[Alternating matrix predicates]
For \(m\in \mathbb{N}\), let
\[
Q_i=
\begin{cases}
\exists,&i\text{ odd},\\
\forall,&i\text{ even}.
\end{cases}
\]
Define
\[
\Xi_m:\Omega_1\to\{0,1\}
\]
by $\Xi_m(A)=1$, if and only if
\[
(Q_1 n_1)(Q_2 n_2)\cdots(Q_m n_m)
\bigl[A(\beta_m(n_1,\ldots,n_m))=1\bigr].
\]
The associated SCI computational problem is
\[
\mathcal P_m^{\mathrm{mat}} := (\Xi_m,\Omega_1,(\{0,1\},d_{\mathrm{disc}}),\Lambda_{\mathrm{mat}}),
\]
where $d_{\mathrm{disc}}$ denotes the discrete metric.
\end{definition}

\subsection{Base Maps Over Coordinate Oracles Are Clopen}

\begin{lemma}[Type-G base maps over coordinate oracles are clopen]\label{lem:coordinate-base-clopen}
Let
\[
\Gamma:\Omega_1\to\{0,1\}
\]
be a type-G general algorithm with respect to the coordinate oracle \(\Lambda_{\mathrm{mat}}\). Then $\Gamma^{-1}(\{1\})$ is clopen in \(\Omega_1\).
\end{lemma}

\begin{proof}
Let \(A\in\Omega_1\). By the definition of a general algorithm, the query policy assigns to \(A\) a finite set
\[
\Lambda_\Gamma(A)\subseteq\Lambda_{\mathrm{mat}}.
\]
Equivalently, there is a finite set of matrix coordinates
\[
D_A\subseteq\mathbb N^2
\]
such that
\[
\Lambda_\Gamma(A)=\{\operatorname{ev}_{a,b}:(a,b)\in D_A\}.
\]
Define the basic cylinder
\[
N_A:= \{B\in\Omega_1:B(a,b)=A(a,b)\text{ for all }(a,b)\in D_A\}.
\]
This set is clopen. By finite-information dependence in the definition of a general algorithm, if \(B\in N_A\), then
\[
\Gamma(B)=\Gamma(A).
\]
Indeed, \(B\) agrees with \(A\) on all queries used by the algorithm at \(A\). Therefore
\[
A\in\Gamma^{-1}(\{1\}) \quad\Longrightarrow\quad N_A\subseteq\Gamma^{-1}(\{1\}),
\]
and
\[
A\notin\Gamma^{-1}(\{1\}) \quad\Longrightarrow\quad N_A\subseteq\Omega_1\setminus\Gamma^{-1}(\{1\}).
\]
Thus both \(\Gamma^{-1}(\{1\})\) and its complement are open. Hence \(\Gamma^{-1}(\{1\})\) is clopen.
\end{proof}

\begin{lemma}[Height \(k\) coordinate towers yield \(\Delta^0_{k+1}\) sets]\label{lem:height-k-delta}
Let
\[
\Theta:\Omega_1\to\{0,1\}
\]
be computed by a type-G tower of height \(k\) over the coordinate oracle \(\Lambda_{\mathrm{mat}}\). Then
\[
\Theta^{-1}(\{1\})\in\Delta^0_{k+1}(\Omega_1).
\]
\end{lemma}

\begin{proof}
We prove this by induction on \(k\). For \(k=0\), the map \(\Theta\) is a single type-G base algorithm. By \cref{lem:coordinate-base-clopen}, its acceptance set is clopen, hence belongs to \(\Delta^0_1\).

Assume the result holds for towers of height \(k-1\). Let \(\Theta\) be computed by a height-\(k\) tower. Then there are maps
\[
\Theta_n:\Omega_1\to\{0,1\}
\]
such that
\[
\Theta(A)=\lim_{n\to\infty}\Theta_n(A)
\]
for $A\in\Omega_1$, and each \(\Theta_n\) is computed by a height-\((k-1)\) tower. By the induction hypothesis,
\[
A_n:=\Theta_n^{-1}(\{1\})\in\Delta^0_k(\Omega_1).
\]
Since the target \(\{0,1\}\) is discrete, pointwise convergence is eventual constant. Therefore
\[
\Theta^{-1}(\{1\}) = \bigcup_{N=1}^\infty \bigcap_{n\ge N} A_n.
\]
Because each \(A_n\in\Delta^0_k\), the right-hand side belongs to \(\Sigma^0_{k+1}\).

Similarly,
\[
\Omega_1\setminus\Theta^{-1}(\{1\}) = \bigcup_{N=1}^\infty \bigcap_{n\ge N}(\Omega_1\setminus A_n),
\]
and each \(\Omega_1\setminus A_n\in\Delta^0_k\). Hence the complement also belongs to \(\Sigma^0_{k+1}\). Thus
\[
\Theta^{-1}(\{1\})\in\Delta^0_{k+1}.
\]
\end{proof}

\subsection{Borel Complexity Of The Alternating Predicates}

\begin{lemma}[\(\Sigma^0_m\)-completeness]\label{lem:Xi-m-complete}
For every \(m\in \mathbb{N}\),
\[
\Xi_m^{-1}(\{1\})
\]
is \(\Sigma^0_m\)-complete in the Cantor space \(\Omega_1\). In particular,
\[
\Xi_m^{-1}(\{1\})\notin\Delta^0_m(\Omega_1).
\]
\end{lemma}

\begin{proof}
First, membership in \(\Sigma^0_m\) follows directly from the definition of \(\Xi_m\). The atomic condition
\[
A(\beta_m(n_1,\ldots,n_m))=1
\]
is clopen. Alternating countable unions and intersections beginning with a countable union give a \(\Sigma^0_m\) set.

We prove hardness using the standard normal form theorem for Borel sets in zero-dimensional Polish spaces (follows from the existence of universal sets for the finite
Borel hierarchy, see, for example, \cite[Theorem~3.6.6]{Sri98}): for every \(\Sigma^0_m\) set \(B\subseteq 2^{\mathbb N}\) there is a clopen predicate
\[
R(y;n_1,\ldots,n_m)
\]
such that
\[
y\in B \quad\Longleftrightarrow\quad (Q_1n_1)\cdots(Q_mn_m)\,R(y;n_1,\ldots,n_m).
\]
Define
\[
\Theta:2^{\mathbb N}\to\Omega_1
\]
coordinatewise by
\[
\Theta(y)(\beta_m(n_1,\ldots,n_m)) :=
\begin{cases}
1,&R(y;n_1,\ldots,n_m),\\
0,&\text{otherwise}.
\end{cases}
\]
Each coordinate of \(\Theta\) is continuous because \(R(\,\cdot\,;n_1,\ldots,n_m)\) is clopen. Hence \(\Theta\) is continuous as a map into the product space \(\Omega_1\). By construction,
\[
y\in B \quad\Longleftrightarrow\quad \Theta(y)\in\Xi_m^{-1}(\{1\}).
\]
Thus every \(\Sigma^0_m\) subset of \(2^{\mathbb N}\) continuously reduces to \(\Xi_m^{-1}(\{1\})\), proving \(\Sigma^0_m\)-hardness. Therefore \(\Xi_m^{-1}(\{1\})\) is \(\Sigma^0_m\)-complete.

The strictness of the finite Borel hierarchy on Cantor space implies that a \(\Sigma^0_m\)-complete set is not in \(\Delta^0_m\).
\end{proof}

\subsection{Exact Type-G Height}

\begin{lemma}[Finite Boolean combinations preserve tower height]\label{lem:boolean-combinations-same-height}
Let \(\Theta_1,\ldots,\Theta_N:\Omega_1\to\{0,1\}\) be decision maps each computed by a type-G tower of height at most \(k\). Then every Boolean combination of \(\Theta_1,\ldots,\Theta_N\) is computed by a type-G tower of height at most \(k\).
\end{lemma}

\begin{proof}
It suffices to treat negation and finite conjunctions. For negation, let \((\Gamma_{n_k,\ldots,n_1})\) be a height-\(k\) tower for \(\Theta\). Then
\[
        1-\Theta(A)
        =
        \lim_{n_k\to\infty}\cdots\lim_{n_1\to\infty}
        \bigl(1-\Gamma_{n_k,\ldots,n_1}(A)\bigr).
\]
For each deepest-level general algorithm \(\Gamma_{n_k,\ldots,n_1}\), the map
\[
        A\mapsto 1-\Gamma_{n_k,\ldots,n_1}(A)
\]
uses the same finite query set and then applies the finite post-processing \(b\mapsto 1-b\). Hence it is again a general algorithm.

For finite conjunctions, first pad all towers to height exactly \(k\) by adding dummy outer indices if necessary. For fixed multi-index \((n_k,\ldots,n_1)\), run the finitely
many deepest-level general algorithms for \(\Theta_1,\ldots,\Theta_N\) at the same
multi-index and query the union of their finite query sets. This union is finite.
Output the product of the obtained \(0/1\)-values. Agreement on the union of the
queried evaluations forces agreement for each constituent algorithm, hence forces
agreement of the product and of the constituent query sets.  Thus the product map is
again a general algorithm.

The iterated limits commute with finite products in the discrete target
\(\{0,1\}\), because convergence in \(\{0,1\}\) is eventual constancy at each limiting
stage. Hence the resulting tower computes the finite conjunction. Disjunctions follow
by De Morgan's law or by replacing product by maximum.
\end{proof}

\begin{theorem}[Exact finite type-G heights of the matrix source problems]\label{thm:Xi-m-exact-height}
For every \(m\in \mathbb{N}\),
\[
\mathrm{SCI}_G(\mathcal P_m^{\mathrm{mat}})=m.
\]
\end{theorem}

\begin{proof}
We prove the upper and lower bounds separately.

\smallskip
\noindent\textbf{Upper bound:}
For a finite tuple \(s=(n_1,\ldots,n_r)\), \(0\le r\le m\), define the tail predicate
\[
\Xi_{m,r}(A;s)
\]
as follows: If \(r=m\), put
\[
\Xi_{m,m}(A;n_1,\ldots,n_m) := A(\beta_m(n_1,\ldots,n_m)).
\]
If \(r<m\), define recursively
\[
\Xi_{m,r}(A;s)=1
\]
if and only if
\[
(Q_{r+1} n_{r+1})\cdots(Q_m n_m) \bigl[A(\beta_m(n_1,\ldots,n_m))=1\bigr].
\]
Thus \(\Xi_m(A)=\Xi_{m,0}(A;\varnothing)\).

We prove by backward induction on \(r\) that \(\Xi_{m,r}(\,\cdot\,;s)\) is computable by a type-G tower of height at most \(m-r\), uniformly in the finite parameter \(s\).

For \(r=m\), the map
\[
A\mapsto A(\beta_m(n_1,\ldots,n_m))
\]
is a single coordinate query, hence a height-\(0\) general algorithm. Assume now \(r<m\). If \(Q_{r+1}=\exists\), then
\[
\Xi_{m,r}(A;s) = \lim_{N\to\infty} \max_{1\le n\le N}\Xi_{m,r+1}(A;s,n).
\]
If \(Q_{r+1}=\forall\), then
\[
\Xi_{m,r}(A;s) = \lim_{N\to\infty} \min_{1\le n\le N}\Xi_{m,r+1}(A;s,n).
\]
For each fixed \(N\), the finite maximum or minimum is a finite Boolean combination of maps computable by towers of height at most \(m-r-1\). By \cref{lem:boolean-combinations-same-height}, it is computable by a tower of height at most \(m-r-1\). The additional limit over \(N\) gives height at most \(m-r\). Taking \(r=0\) gives a height-\(m\) tower for \(\Xi_m\).

\smallskip
\noindent\textbf{Lower bound:}
Suppose, toward a contradiction, that
\[
\mathrm{SCI}_G(\mathcal P_m^{\mathrm{mat}})\le m-1.
\]
Then \(\Xi_m\) is computed by a type-G tower of height at most \(m-1\). By \cref{lem:height-k-delta},
\[
\Xi_m^{-1}(\{1\})\in\Delta^0_m(\Omega_1).
\]
This contradicts \cref{lem:Xi-m-complete}, which says that \(\Xi_m^{-1}(\{1\})\) is \(\Sigma^0_m\)-complete and therefore not in \(\Delta^0_m\).
Thus
\[
\mathrm{SCI}_G(\mathcal P_m^{\mathrm{mat}})\ge m.
\]
Combining both inequalities proves the theorem.
\end{proof}

\begin{definition}[Tagged union source problem]
Let
\[
\Omega_\infty:=\{(m,A) : m\in\mathbb N,\ A\in\Omega_1\}.
\]
Equip \(\Omega_\infty\) with the oracle consisting of the tag map
\[
\tau(m,A):=m
\]
together with all matrix-coordinate evaluations
\[
        \operatorname{ev}_{a,b}^{\infty}:\Omega_\infty\to\{0,1\}, \qquad
        \operatorname{ev}_{a,b}^{\infty}(m,A):=A(a,b).
\]

Set
\[
        \Lambda_\infty := \{\tau\}\cup \{\operatorname{ev}_{a,b}^{\infty} : (a,b)\in\mathbb N^2\}.
\]
Define
\[
\Xi_\infty(m,A):=\Xi_m(A).
\]
Let
\[
\mathcal P_\infty^{\mathrm{mat}} := (\Xi_\infty,\Omega_\infty,(\{0,1\},d_{\mathrm{disc}}),\Lambda_\infty)
\]
be the corresponding SCI computational problem.
\end{definition}

\begin{theorem}[Unbounded finite heights give infinite type-G height]\label{thm:Xi-infty-height}
One has
\[
\mathrm{SCI}_G(\mathcal P_\infty^{\mathrm{mat}})=\infty.
\]
\end{theorem}

\begin{proof}
Assume for contradiction that
\[
        \SCIG(\mathcal P_\infty^{\mathrm{mat}})=k<\infty.
\]
Restrict a height-\(k\) tower for \(\mathcal P_\infty^{\mathrm{mat}}\) to the slice
\[
        \Omega_{k+1}:=\{(k+1,A) : A\in\Omega_1\}.
\]
On this slice, the tag query \(\tau\) is constant with value \(k+1\), and every matrix
coordinate query \(\operatorname{ev}_{a,b}\) is exactly the corresponding coordinate
evaluation of \(A\). Therefore each deepest-level general algorithm in the restricted
tower is simulated by a deepest-level general algorithm over the coordinate oracle
\(\Lambda_{\mathrm{mat}}\): tag queries are answered by the fixed constant \(k+1\), and
coordinate queries are passed to the corresponding coordinate oracle on \(A\).

Thus the restricted tower is a height-\(k\) type-\(G\) tower computing \(\Xi_{k+1}\) on
\(\Omega_1\). This contradicts \cref{thm:Xi-m-exact-height}, which gives
\[
        \SCIG(\mathcal P_{k+1}^{\mathrm{mat}})=k+1.
\]
Hence no finite-height type-\(G\) tower computes \(\mathcal P_\infty^{\mathrm{mat}}\), and
\[
        \SCIG(\mathcal P_\infty^{\mathrm{mat}})=\infty.
\]
\end{proof}

\section{Uniformity, Minimality, And Pure \texorpdfstring{\(\mathcal R\)-\(\Lim\)}{R-Lim} Implemented SCI}\label{sec:uniformity-minimality}

\subsection{Representations, Atoms, And Regularity Classes}\label{sec:representations-atoms-regularity}
To answer the proposed question, we have to define first, what we actually mean by an computability model in general.

\begin{definition}[Computability model]\label{def:compmodel}
Fix a class of \textit{types} $\mathcal{T}$ and for each $X,Y\in\mathcal{T}$ a class $\mathbb{M}(X,Y)$
of morphisms (partial or multi-valued if desired).
We call $\mathbb{M}$ a \textit{computability model} if
\begin{enumerate}
\item (\textbf{Well-definedness}) Each $\mathbb{M}(X,Y)$ is a \textit{definite} set of partial or multivalued maps from $X$ to $Y$.
\item (\textbf{Identity}) $\id_X\in\mathbb{M}(X,X)$ for all $X$.
\item (\textbf{Composition}) If $f\in\mathbb{M}(X,Y)$ and $g\in\mathbb{M}(Y,Z)$ then $g\circ f\in\mathbb{M}(X,Z)$.
\item (\textbf{Basic data operations}) Pairing and projections for products are in $\mathbb{M}$ (at least at the level of names),
so that one can feed outputs of one procedure as inputs of another.
\end{enumerate}
\end{definition}

To have everything compact here, we also recall again, what a representation and realizer is.
\begin{definition}[Represented space and realizer]
A \textit{represented space} is $X=(|X|,\delta_X)$ where $\delta_X:\subseteq\Baire\to |X|$ is a partial surjection.
A partial function $F:\subseteq\Baire\to\Baire$ \textit{realizes} \(f:\subseteq |X|\rightrightarrows |Y|\), written \(F\vdash f\), if for every
\(p\in\operatorname{dom}(\delta_X)\) with \(\delta_X(p)\in\operatorname{dom}(f)\), one has \(F(p)\in\operatorname{dom}(\delta_Y)\) and
\[
        \delta_Y(F(p))\in f(\delta_X(p)).
\]
\end{definition}

\begin{definition}[Representation-induced topology]
Given $(X,\delta_X)$, define a topology $\tau_{\delta_X}$ on $|X|$ by declaring $U\subseteq |X|$ open iff
$\delta_X^{-1}(U)\subseteq\Baire$ is open (in the product topology on $\Baire$).
\end{definition}

Now we first recall a well known but for our purpose very crucial lemma.
\begin{lemma}[Type-2 computable implies being continuous]\label{lem:comp-implies-cont}
Let $f:|X|\to |Y|$ be Type-2 computable between represented spaces.
Then $f:(|X|,\tau_{\delta_X})\to(|Y|,\tau_{\delta_Y})$ is continuous.
\end{lemma}
\begin{proof}
Let $U\subseteq |Y|$ be open. Then $\delta_Y^{-1}(U)$ is open in $\Baire$.
Let $F\vdash f$ be a computable realizer; any computable $F$ is continuous on $\Baire$.
Hence $F^{-1}(\delta_Y^{-1}(U))$ is open in $\Baire$.
But $F^{-1}(\delta_Y^{-1}(U))=\delta_X^{-1}(f^{-1}(U))$ by the realizer condition.
Thus $\delta_X^{-1}(f^{-1}(U))$ is open, i.e.\ $f^{-1}(U)$ is open in $\tau_{\delta_X}$.
\end{proof}

By this lemma, we will now show, that the $\mathrm{SCI}_{\mathrm A}$ necessarily has to fix a representation.
\begin{theorem}[Representation dependence: \textit{any} function can be made computable]\label{thm:any-f-computable-by-rep}
Let $Y=(|Y|,\delta_Y)$ be any represented space and let $f:|X|\to |Y|$ be any function on a bare set $|X|$.
Assume $|X|$ has \textit{some} representation $\delta_X:\subseteq\Baire\to |X|$ (not necessarily admissible).
Then there exists another representation $\delta_X^f:\subseteq\Baire\to |X|$ such that $f:(|X|,\delta_X^f)\to(|Y|,\delta_Y)$ is Type-2 computable.
\end{theorem}
\begin{proof}
Fix a computable pairing $\pair:\Baire\times\Baire\to\Baire$ with computable projections $\pi_0,\pi_1$.
Define $\delta_X^f:\subseteq\Baire\to |X|$ by
\[
\delta_X^f(\pair(p,q)) := \delta_X(p)
\quad\text{provided that }\delta_Y(q)= f(\delta_X(p)).
\]
This is a partial surjection because for each $x\in|X|$ choose any name $p$ with $\delta_X(p)=x$
and any name $q$ with $\delta_Y(q)=f(x)$, then $\delta_X^f(\pair(p,q))=x$.

Now define $F:\subseteq\Baire\to\Baire$ by $F(r):=\pi_1(r)$.
Then for any $r=\pair(p,q)\in\dom(\delta_X^f)$ we have $\delta_Y(F(r))=\delta_Y(q)=f(\delta_X(p))=f(\delta_X^f(r))$.
Hence $F\vdash f$, and $F$ is computable. Thus $f$ is computable w.r.t.\ $\delta_X^f$.
\end{proof}

\begin{corollary}[Logical necessity of fixing representations]\label{cor:need-representations}
Any notion of \enquote{computable function on $|X|$} that does \textit{not} fix a representation (or at least an equivalence
class of representations) is ill-posed: by \cref{thm:any-f-computable-by-rep} \textit{every} function can be made
computable by a change of representation.
Therefore, specifying representations of inputs and outputs is a logically necessary premise for calling a framework a computability model in the Type-2 sense.
\end{corollary}

\begin{remark}[This is the first minimality premise]
\Cref{cor:need-representations} is a purely formal obstruction: without representations there is no canonical
computability notion, hence no model in the sense of \cref{def:compmodel}.
This is why any comparison of SCI-$A$ to Weihrauch theory \textit{must} first pass to represented spaces
(e.g.\ via the countable evaluation-table representation).
\end{remark}

\begin{definition}[Atomic real interface (BSS-style primitive)]
An \textit{atomic real interface} is an interface in which an algorithm may apply exact predicates
such as $(x<0)$ or $(x=0)$ to real inputs as a single primitive operation.
\end{definition}

The following proposition is a very classical result from computability theory.

\begin{proposition}[Atomic comparison violates Type-2 continuity constraints]\label{prop:atomic-breaks}
Let $\mathrm{sgn}:\R\to\{-1,0,1\}$ be the sign function with discrete codomain.
Then
\begin{enumerate}
\item $\mathrm{sgn}$ is computable in any model that has an atomic real interface (one comparison to $0$).
\item $\mathrm{sgn}$ is not Type-2 computable w.r.t.\ the standard Cauchy representation of $\R$.
\end{enumerate}
\end{proposition}

\begin{proof}
\textbf{(1):} Immediate from the availability of an exact test $x<0$, $x=0$, $x>0$.

\textbf{(2):} By \cref{lem:comp-implies-cont}, any Type-2 computable map $\R\to\{-1,0,1\}$ (discrete)
must be continuous. But $\mathrm{sgn}$ is discontinuous at $0$, hence not Type-2 computable.
\end{proof}

\begin{corollary}[Second minimality premise: to compare to Weihrauch, reals must be given by names]\label{cor:names-not-atoms}
Any SCI-$A$ interpretation that treats evaluation outputs as \textit{atomic} reals/complex numbers with exact comparisons
cannot be equivalent to the Type-2/Weihrauch interpretation: it computes discontinuous functions at height $0$
(cf. \cref{prop:atomic-breaks}) that are not even continuous, hence not Type-2 computable.
Therefore, \enquote{oracle returns names/approximations} (not atoms) is logically necessary for Weihrauch comparability.
\end{corollary}

\begin{remark}[Connection to SCI$_A$ vs. BSS/TTE]
This is exactly the semantic difference behind the BSS vs. TTE gap at low heights firstly recognized in \cite[Obs.~43, Thm.~44]{neumann2018topological}. \Cref{cor:names-not-atoms} isolates the \textit{pure logical} reason the gap can exist at all.
\end{remark}

\begin{definition}[Adequate regularity class on $\Baire$]\label{def:adequate-R}
A class $\mathcal R$ of partial functions $\Baire\to\Baire$ is called \textit{adequate} if
\begin{enumerate}
\item it contains all computable pairing and projection maps on $\Baire$;
\item it is closed under composition;
\item it contains all constant maps.
\end{enumerate}
Examples: $\Comp$ (computable), $\Cont$ (continuous), $\Bor$ (Borel measurable).
\end{definition}

Additionally to these three classical examples for $\mathcal R$, we also give an interesting fourth one.
\begin{definition}[Finite-mind-change / limit-computable functionals]\label{def:FMC}
A partial functional $F:\subseteq\NN^\NN\to\NN^\NN$ is called \textit{finite-mind-change computable}
(or \textit{limit computable}) if there exists a \textit{computable} functional
$\Phi:\subseteq\NN^\NN\times\NN\to\NN^\NN$ such that for every $p\in\dom(F)$ and every $n\in\NN$
the sequence $\bigl(\Phi(p,s)(n)\bigr)_{s\in\NN}$ is eventually constant and
\[
F(p)(n)=\lim_{s\to\infty}\Phi(p,s)(n).
\]
Let $\mathsf{FMC}$ denote the class of all finite-mind-change computable functionals
$\subseteq\NN^\NN\to\NN^\NN$.
\end{definition}

\begin{lemma}\label{lem:FMC-adequate}
The class $\mathsf{FMC}$ is adequate in the sense of \cref{def:adequate-R}.
\end{lemma}
\begin{proof}
Pairing/projections and constant maps are computable, hence in $\mathsf{FMC}$ by taking
$\Phi(p,s)$ constant in $s$.
For closure under composition, we use the standard characterization of limit computable
functionals: $F$ is finite-mind-change computable iff $F$ is computable relative to $0'$ by Shoenfield's limit lemma \cite{Sho59}. Since $0'$-computable functionals are closed under composition, $\mathsf{FMC}$ is closed under composition.
\end{proof}

We also recall from the intermediate hierarchy
\begin{definition}[$\mathcal R$-realizable maps]\label{def:R-realizable}
A \textit{represented space} is a pair $X=(|X|,\delta_X)$ where
$\delta_X:\subseteq \Baire\to |X|$ is a partial surjection.
Let $X=(|X|,\delta_X)$ and $Y=(|Y|,\delta_Y)$ be represented spaces and let
$f:\subseteq |X|\to |Y|$ be single-valued.
A partial function $F:\subseteq\Baire\to\Baire$ \textit{realizes} $f$
(written $F \vdash f$) if for all $p\in\dom(\delta_X)$ with $\delta_X(p)\in\dom(f)$,
we have $F(p)\in\dom(\delta_Y)$ and
\[
\delta_Y(F(p)) \;=\; f(\delta_X(p)).
\]
The map $f$ is \textit{$\mathcal R$-realizable} if it has some realizer $F\vdash f$
with $F\in \mathcal R$.
\end{definition}

\begin{lemma}[If $\mathcal{R}$ is not adequate then reduction preorders can fail]\label{lem:R-needed}
If $\mathcal{R}$ omits pairing/projections or is not closed under composition, then the corresponding
$\mathcal{R}$-Weihrauch reducibility (defined by requiring reduction functionals in $\mathcal{R}$)
need not be reflexive or transitive.
\end{lemma}
\begin{proof}
Reflexivity requires projection witnesses (to output the oracle answer unchanged).
Transitivity requires closure of the witness class under composition and pairing.
\end{proof}

\begin{definition}[$\mathcal R$-Weihrauch $\Lim$-rank]\label{def:R-lim-rank}
For any problem $f$ between represented spaces define
\[
\mathrm{rank}_{\mathcal R}(f)
\;:=\;
\min\bigl\{k\in\NN_0:\; f \le_{\mathrm W}^{\mathcal R} \Limi{k}\bigr\},
\]
with value $\infty$ if the set is empty.
\end{definition}

\begin{theorem}[Third minimality premise: unrestricted base post-processing collapses computability]\label{thm:unrestricted-collapses}
Fix any nontrivial set $X$ and consider the SCI problem $\mathcal{P}=(\Xi,\Omega,(\mathcal{M},d),\Lambda)$ with
$\Omega=X$, $\Lambda=\{\id_X\}$ (encoded into $\C$ if desired), and arbitrary target $\Xi:X\to\mathcal{M}$.
If base post-processing is unrestricted (type-$G$ style), then $\mathrm{SCI}_G(\mathcal{P})=0$.
\end{theorem}
\begin{proof}
Define the base algorithm
\[
\Gamma(x):=\Xi(x)
\]
and use the constant one-query policy
\[
\Lambda_\Gamma(x):=\{\mathrm{id}_X\}.
\]
If \(y\in X\) agrees with \(x\) on this query, then
\[
\mathrm{id}_X(y)=\mathrm{id}_X(x),
\]
hence \(y=x\).  Therefore
\[
\Gamma(y)=\Gamma(x)
\]
and also
\[
\Lambda_\Gamma(y)=\Lambda_\Gamma(x).
\]
Thus \((\Gamma,\Lambda_\Gamma)\) is a general algorithm. Since \(\Gamma=\Xi\), this is a height-\(0\) type-G tower.
\end{proof}

\begin{corollary}[Necessity of a restriction class]
To obtain a nontrivial computability model (and to compare to Weihrauch-style reducibilities),
one must restrict base post-processing to a distinguished class $\mathcal{R}$
(e.g.\ computable/continuous/Borel). Otherwise the theory collapses at height $0$
(cf. \cref{thm:unrestricted-collapses}).
\end{corollary}

\subsection{Nonuniform Deepest-Only Towers And The All-Level Over-Repair}\label{sec:nonuniformity-all-level}
This section separates three notions which should not be conflated.

\begin{enumerate}
\item The \textit{standard SCI tower} is the deepest-level tower of \cref{def:standard-tower}. Only the deepest approximants are required to be finite-information algorithms; intermediate maps are semantic pointwise limits.

\item The \textit{all-level finite-information repair} requires every intermediate semantic limit map to be a general algorithm. This repair is too strong: \Cref{prop:all-level-collapse} shows that it collapses every finite positive height to height \(1\).

\item The \textit{pure $\mathcal R -\Lim$ normal form} requires one uniform represented-space preprocessor
\[
        K\in \mathcal R
\]
with
\[
        \Limi{k} \circ K\preceq \mathcal N_{\widehat{\Xi}}.
\]
This is the correct condition for Weihrauch comparability.
\end{enumerate}

The point of this work is therefore not that the standard SCI tower should be replaced by the all-level repair. Rather, the standard SCI tower should be read as a raw extensional
limit skeleton; to obtain a computability model, one must impose uniform implementation.

\begin{proposition}[The all-level repair is strictly stronger than the standard tower]\label{prop:all-level-strictly-stronger}
For \(k\ge2\), the all-level finite-information repair is strictly stronger than the standard deepest-level tower. Already for \(k=2\), there exists a standard height-\(2\) tower whose intermediate layer is not a general algorithm.
\end{proposition}
\begin{proof}
	It suffices to treat the case $k=2$.
	Let
	\[
	X
	:=
	\Bigl\{
	x=(x_{n,m})_{n,m\in\mathbb N}\in\{0,1\}^{\mathbb N^2}
	:
	\lim_{n\to\infty}\sum_{m=1}^{\infty}2^{-m}x_{n,m}\text{ exists}
	\Bigr\},
	\]
	let $(\mathcal M,d)=([0,1],|\cdot|)$, and let
	\[
	\Lambda:=\{\pi_{n,m}:X\to\{0,1\}\mid \pi_{n,m}(x)=x_{n,m}\}.
	\]
	Define
	\[
	\Gamma_n(x):=\sum_{m=1}^{\infty}2^{-m}x_{n,m},
	\qquad
	\Gamma_{n,k}(x):=\sum_{m=1}^{k}2^{-m}x_{n,m},
	\qquad
	\Xi(x):=\lim_{n\to\infty}\Gamma_n(x).
	\]
	The target \(\Xi\) is well-defined by the definition of \(X\), and the consistency condition holds because the coordinate evaluations in \(\Lambda\) separate points of \(X\).
	For every fixed pair $(n,k)$, the map $\Gamma_{n,k}$ depends only on the finitely many
	coordinates $\pi_{n,1},\ldots,\pi_{n,k}$ and is obtained from them by finitely many
	arithmetic operations.
	Hence $(\Gamma_n,\Gamma_{n,k})$ is a height-$2$ tower in the sense of \cite[Definition~3.1]{hansen2011solvability}.
	
	However, $\Gamma_n$ is not a general algorithm in the sense of \cref{def:gen-alg}.
	Indeed, suppose $(\Gamma_n,\Lambda_{\Gamma_n})$ were such a general algorithm.
	Fix $x\equiv 0$.
	Then $\Lambda_{\Gamma_n}(x)\subseteq \Lambda$ is finite, so there exists $m_0\in\mathbb N$
	with $\pi_{n,m_0}\notin \Lambda_{\Gamma_n}(x)$.
	Let $y\in X$ be given by
	\[
	y_{n,m_0}=1,
	\qquad
	y_{r,s}=0 \quad \text{for all }(r,s)\neq (n,m_0).
	\]
	Then
	\[
	\pi(y)=\pi(x)\qquad\text{for all }\pi\in\Lambda_{\Gamma_n}(x),
	\]
	but
	\[
	\Gamma_n(y)-\Gamma_n(x)=2^{-m_0}\neq 0.
	\]
	This contradicts finite-information dependence in \cref{def:gen-alg}.
	Hence \(\Gamma_n\) is not a general algorithm.  Thus the tower is admissible as a
standard deepest-level height-\(2\) tower, but not as an all-level finite-information
tower.
	
	For larger $k$, one simply keeps all additional lower indices dummy.
\end{proof}

The previous proposition shows that the all-level repair is too strong. We now show
the opposite problem: a represented deepest-only tower with individually computable
deepest estimators is still too weak to imply a Weihrauch normal form.  The missing
ingredient is uniformity of the whole indexed approximation table.

\begin{definition}[Represented weak Hansen package]\label{def:weak-Hansen-represented}
	Fix a countable-$\Lambda$ SCI computational problem
	\[
	\mathcal P= \bigl( \Xi,\Omega,(\mathcal M,d),\Lambda \bigr)
	\]
	with fixed represented information space $(I_\Lambda,\delta_{I_\Lambda})$ and a fixed
	representation $\delta_{\mathcal M}$ of $\mathcal M$, as in \cref{def:SCIAR-package}.
	Let $\mathcal R\subseteq \mathbb N^{\mathbb N}\to\mathbb N^{\mathbb N}$ be a class of realizers.
	A \textit{represented weak Hansen $\mathcal R$-tower of height $k\in\mathbb N$} means a family of
	maps
	\[
	\Gamma_{n_k},\ \Gamma_{n_k,n_{k-1}},\ \dots,\ \Gamma_{n_k,\dots,n_1}\colon \Omega\to \mathcal M
    \]
    for $n_1,\dots,n_k\in\mathbb N$	such that
	\begin{enumerate}
		\item the same pointwise iterated-limit equalities as in \cref{def:standard-tower} hold, i.e.
		\[
		\Xi(A)=\lim_{n_k\to\infty}\Gamma_{n_k}(A),\qquad
		\Gamma_{n_k}(A)=\lim_{n_{k-1}\to\infty}\Gamma_{n_k,n_{k-1}}(A),
		\]
		and so on down to depth $k$;
		\item only the deepest maps $\Gamma_{n_k,\dots,n_1}$ are required to be
		$\mathcal R$-realizable general algorithms (equivalently: after fixing names, each deepest map has an $\mathcal R$-realizer and satisfies \cref{def:gen-alg}); no condition is imposed on the intermediate maps
		\[
		\Gamma_{n_k},\ \Gamma_{n_k,n_{k-1}},\ \dots,\ \Gamma_{n_k,\dots,n_2}.
		\]
	\end{enumerate}
	This is the literal represented-space repair of \cite[Definition~3.1]{hansen2011solvability} obtained by keeping Hansen's deepest-estimator clause and fixing the extra represented-space data: representations, names, and a class \(\mathcal R\) of admissible realizers.
\end{definition}

\begin{proposition}[The deepest-only Hansen clause does not imply a pure $\mathcal R$-$\Lim$ witness]\label{prop:weak-Hansen-no-K}
	Fix $\mathcal R =\mathrm{Comp}$.  There exists a countable-$\Lambda$ SCI computational problem
	\[
	\mathcal P=\bigl( \Xi,\Omega,(\mathcal M,d),\Lambda \bigr)
	\]
	with fixed represented information space and fixed represented output space such that
	\begin{enumerate}
		\item $\mathcal P$ admits a represented weak Hansen $\mathrm{Comp}$-tower of height $2$ in the sense of \cref{def:weak-Hansen-represented};
		\item every deepest estimator $\Gamma_{n,m}$ is in fact a fixed-query computable general
		algorithm;
		\item every intermediate estimator $\Gamma_n$ fails to be a general algorithm in the
		sense of \cref{def:gen-alg};
		\item nevertheless there is \textit{no} computable function
		\[
		K\colon\subseteq \mathbb N^{\mathbb N}\to\mathbb N^{\mathbb N}
		\]
		with
		\[
		\Limi{2} \circ K \preceq \mathcal N_{\widehat{\Xi}}
		\]
		in the sense of \cref{def:SCIAR-package}.
	\end{enumerate}
	In plain words, even after fixing the represented-space data required in \cref{sec:uniformity-minimality},
	Hansen's deepest estimator clause is still too weak to yield Weihrauch comparability. It does \textit{not} imply the existence of the uniform preprocessor $K$ required in \cref{def:SCIAR-package} and \cref{thm:equivalence}.
\end{proposition}

\begin{proof}
	We build and verify a counterexample explicitly in six steps.
	
	\smallskip
	\noindent
	\textbf{Step~1: the represented input space:}
	Let
	\[
	\Omega:=\Bigl\{x\in 2^{\mathbb N\times\mathbb N_0} : (\forall j\in\mathbb N_0)\ \lim_{n\to\infty}x(n,j)
	\text{ exists in }\{0,1\}\Bigr\}.
	\]
	Equip $\Omega$ with the subspace representation inherited from the Cantor space
	$2^{\mathbb N\times\mathbb N_0}$.
	Let
	\[
	\Lambda:=\{\pi_{n,j}:\Omega\to\{0,1\}\subseteq\mathbb C \mid n\in\mathbb{N},\ j\in\mathbb{N}_0 \},
	\qquad \pi_{n,j}(x):=x(n,j).
	\]
	Thus the induced represented information space is the standard name space for the
	countable table of bits of $x$. Via this coordinate interface, the factor space $I_\Lambda$ identifies canonically with $\Omega$ itself; in the sequel we silently use this identification when writing $\widehat\Xi$.
	
	\smallskip
	\noindent
	\textbf{Step~2: the represented output space:}
	Let $\mathcal C:=2^{\mathbb N_0}$ be the Cantor space with its standard representation and let
	\[
	\mathcal M:=\mathcal C\times \mathcal C
	\]
	with the usual product representation and any compatible product metric.
	Write $0^{\omega}\in\mathcal C$ for the all-zero sequence.
	
	Now consider the set
	\[
	S:=\Bigl\{u\in\mathcal C : \exists q\in\operatorname{dom}(\Limi{2})\cap\mathbb N^{\mathbb N}
	\text{ computable with }\delta_{\mathcal M}(\Limi{2}(q))=(u,v)
	\text{ for some }v\in\mathcal C\Bigr\}.
	\]
	There are only countably many computable $q$, hence $S$ is countable.
	Choose once and for all some
	\[
	p\in \mathcal C\setminus S.
	\]
	For each $n\in\mathbb N$ define the finite-prefix approximation
	\[
	p_n(j):=
	\begin{cases}
		p(j), & j<n,\\
		0, & j\ge n,
	\end{cases}
	\qquad (j\in\mathbb N_0),
	\]
	and for $n,m\in\mathbb N$ define
	\[
	p_{n,m}:=p_{\min\{n,m\}}.
	\]
	Then each $p_n$ and each $p_{n,m}$ is a computable point of Cantor space (its first
	finitely many bits are hard-coded into the program, and all later bits are $0$), and
	\[
	\lim_{m\to\infty}p_{n,m}=p_n,
	\qquad
	\lim_{n\to\infty}p_n=p
	\]
	in the Cantor space.
	
	\smallskip
	\noindent
	\textbf{Step~3: the target map and the weak Hansen family:}
	Define
	\[
	L\colon \Omega\to \mathcal C,
	\qquad
	L(x)(j):=\lim_{n\to\infty}x(n,j) \qquad (j\in\mathbb{N}_0),
	\]
	which is well-defined by the definition of $\Omega$. For every $n,m\in\mathbb N$ define
	\[
	R_n(x)(j):=x(n,j),
	\qquad
	R_{n,m}(x)(j):=
	\begin{cases}
		x(n,j), & j<m,\\
		0, & j\ge m.
	\end{cases}
	\]
	Now set
	\[
	\Gamma_n(x):=(p_n,R_n(x)),
	\qquad
	\Gamma_{n,m}(x):=(p_{n,m},R_{n,m}(x)),
	\qquad
	\Xi(x):=(p,L(x)).
	\]
	We claim that these maps form a represented weak Hansen $\mathrm{Comp}$-tower of
	height $2$.
	
	First, for each fixed $n$ and $x\in\Omega$ we have
	\[
	\lim_{m\to\infty}\Gamma_{n,m}(x)=\Gamma_n(x).
	\]
	Indeed, the first coordinate satisfies $p_{n,m}\to p_n$ and the second coordinate
	satisfies $R_{n,m}(x)\to R_n(x)$, because the truncations of a binary sequence converge to
	that sequence in the product metric on the Cantor space.
	Second, for each $x\in\Omega$ we have
	\[
	\lim_{n\to\infty}\Gamma_n(x)=\Xi(x).
	\]
	Again the first coordinate satisfies $p_n\to p$, while the second coordinate satisfies
	$R_n(x)\to L(x)$ since for every fixed coordinate $j$ the sequence $n\mapsto x(n,j)$ is
	eventually constant and therefore the convergence is coordinatewise on every finite
	prefix, hence in the Cantor metric.
	Thus the required nested limit equalities hold.
	
	\smallskip
	\noindent
	\textbf{Step~4: the deepest estimators are fixed-query computable general algorithms:}
	Fix $n,m$.  The map $\Gamma_{n,m}$ depends only on the finitely many coordinates
	\[
	\pi_{n,0},\pi_{n,1},\dots,\pi_{n,m-1}.
	\]
	Indeed, the first component $p_{n,m}$ is constant, and the second component
	$R_{n,m}(x)$ only uses the first $m$ entries of the $n$-th row of $x$.
	Hence $\Gamma_{n,m}$ is a fixed-query general algorithm in the sense of
	\cref{def:gen-alg}, witnessed by the constant query policy
	\[
	\Lambda_{\Gamma_{n,m}}(x):=\{\pi_{n,0},\dots,\pi_{n,m-1}\}.
	\]
	Moreover $\Gamma_{n,m}$ is computable: from a name of $x$ one can compute the queried
	bits $x(n,0),\dots,x(n,m-1)$, write down the computable constant $p_{n,m}$, and output a
	name of the pair $(p_{n,m},R_{n,m}(x))$.
	Therefore every deepest estimator is a fixed-query computable general algorithm.
	
	\smallskip
	\noindent
	\textbf{Step~5: the intermediate estimators are not general algorithms:}
	Fix $n\in\mathbb N$. We show that $\Gamma_n$ is not a general algorithm.
	Suppose, towards a contradiction, that some finite query policy
	$\Lambda_{\Gamma_n}(x)$ witnesses \cref{def:gen-alg} for $\Gamma_n$ on some input
	$x\in\Omega$.
	Because $\Lambda_{\Gamma_n}(x)$ is finite, there exists $j_0\in\mathbb N$ such that the
	coordinate query $\pi_{n,j_0}$ does not belong to $\Lambda_{\Gamma_n}(x)$.
	Define $y\in\Omega$ by flipping exactly that one bit in the $n$-th row, i.e.
	\[
	y(i,j):=
	\begin{cases}
		1-x(n,j_0), & (i,j)=(n,j_0),\\
		x(i,j), & \text{otherwise.}
	\end{cases}
	\]
	Then $x$ and $y$ agree on all queries in $\Lambda_{\Gamma_n}(x)$, but
	\[
	\Gamma_n(x)=(p_n,R_n(x))\neq (p_n,R_n(y))=\Gamma_n(y)
	\]
	because the second coordinate differs at position $j_0$.
	This contradicts the finite-information dependence clause in \cref{def:gen-alg}.
	Hence $\Gamma_n$ is not a general algorithm. This contradicts the finite-information dependence clause in
\cref{def:gen-alg}. Hence \(\Gamma_n\) is not a general algorithm.
Since \(n\) was arbitrary, the intermediate semantic layers of this tower are not general
algorithms. Thus the tower is admissible as a standard deepest-level tower, but it is not
admissible as an all-level finite-information tower.
	
	\smallskip
	\noindent
	\textbf{Step~6: no computable $K$ as in \cref{def:SCIAR-package} exists:}
Assume for contradiction that there exists a computable functional
\[
K\colon\subseteq \mathbb N^{\mathbb N}\to\mathbb N^{\mathbb N}
\]
such that
\[
\Limi{2}\circ K \preceq \mathcal N_{\widehat{\Xi}}.
\]
Let $x^{(0)}\in\Omega$ be the computable zero matrix, and set
\[
\widehat{x}^{(0)}:= \mathrm{Ev}_{\Lambda}(x^{(0)})\in I_{\Lambda}.
\]
Let
\[
r^{(0)}\in\dom(\delta_{I_{\Lambda}})
\]
be a computable name of $\widehat{x}^{(0)}$, i.e.
\[
\delta_{I_{\Lambda}}(r^{(0)})=\widehat{x}^{(0)}.
\]
Since $K$ is computable, the sequence
\[
q:=K(r^{(0)})
\]
is computable. By the equivalent formulation in \cref{def:SCIAR-package}, the assumption
\[
\Limi{2}\circ K \preceq \mathcal N_{\widehat{\Xi}}
\]
implies that $q\in\dom(\Limi{2})$ and
\[
\delta_{\mathcal M}(\Limi{2}(q))=\widehat{\Xi}(\widehat{x}^{(0)}).
\]
By the canonical identification $I_{\Lambda}\cong\Omega$ from Step~1,
\[
\widehat{\Xi}(\widehat{x}^{(0)})=\Xi(x^{(0)})=(p,0^\omega).
\]
Hence
\[
\delta_{\mathcal M}(\Limi{2}(q))=(p,0^\omega).
\]
Therefore, by the definition of $S$ in Step~2, we have $p\in S$, contradicting
\[
p\in\mathcal C\setminus S.
\]
This contradiction shows that no such computable $K$ exists.
\end{proof}

\subsection{Failure Of The Three Naive Effective Readings}\label{sec:three-naive-readings}

\begin{theorem}[Failure of the three naive effective SCI readings]\label{thm:logical-necessity-uniform-implementation}
Suppose an SCI formalism is intended to recover Weihrauch/Type-2 computation as a special case in the following minimal sense:
\begin{enumerate}
\item countable-evaluation SCI computational problems can be converted into represented information problems
\[
        \widehat{\Xi}: I_\Lambda\to \mathcal M;
\]
\item finite effective SCI height implies finite rank along the Weihrauch iterated-limit chain;
\item the formalism still allows nontrivial exact finite heights \(2,3,\ldots\).
\end{enumerate}
Then none of the following three structures is sufficient as the main effective
definition:
\begin{enumerate}
\item unrestricted raw type-G SCI;
\item represented deepest-only towers with merely individually computable deepest estimators;
\item all-level finite-query exact-output towers.
\end{enumerate}
\end{theorem}

\begin{proof}
Unrestricted raw type-G is too weak: by \cref{thm:finite-factor-scig0}, any target map factoring through a finite transcript is computable at raw type-G height \(0\), regardless of the descriptive or Type-2 complexity of the finite-transcript post-processor. \Cref{thm:scig0-but-tteinf} gives an explicit analytic non-Borel decision example with raw \(\SCIG=0\) but infinite Weihrauch-SCI rank.

Deepest-only represented towers with merely individually computable deepest estimators are still too weak: \Cref{prop:weak-Hansen-no-K} constructs a represented weak Hansen computable tower of height \(2\) whose deepest estimators are computable, but for which no computable
\[
        K:\subseteq\mathbb N^{\mathbb N}\to\mathbb N^{\mathbb N}
\]
satisfies
\[
        \Limi{2} \circ K\preceq \mathcal N_{\widehat{\Xi}}.
\]
Thus finite deepest-only computable height does not imply finite Weihrauch rank.

All-level finite-query exact-output towers are too strong: by \cref{prop:all-level-collapse}, every finite positive all-level height collapses to height \(1\). Hence this repair cannot preserve exact finite heights \(2,3,\ldots\).

Thus the first reading is too weak, the second reading is still non-uniform, and the third reading is too strong. The remaining repair used in this paper is to require one admissible procedure producing the whole indexed approximation table. In the Weihrauch case this is expressed by the existence of a single admissible preprocessor \(K\) with
\[
        \Limi{k} \circ K\preceq \mathcal N_{\widehat{\Xi}}.
\]
\end{proof}

The conclusion is not that these notions are useless. Raw type-\(G\) is useful as an
information-theoretic skeleton, and all-level towers are useful as a diagnostic
strengthening. The point is that neither gives the right effective Weihrauch-comparable
semantics by itself.

\subsection{Pure \texorpdfstring{\(\mathcal R\)-\(\Lim\)}{R-Lim} Implemented SCI And The Bridge Theorem}\label{sec:pure-R-lim}
The next definition isolates the name-level tower notion used for Weihrauch comparability.
Unlike a raw SCI tower, it requires one admissible preprocessor producing the whole input
to the iterated limit operator. Thus the bridge theorem below is not merely a rephrasing of raw \(\SCIA\); it identifies the additional uniformity needed for a represented-space computation model.

\begin{definition}[Pure $\mathcal R-\Lim$ implemented SCI package]\label{def:SCIAR-package}
Let $\mathcal{P}=\bigl( \Xi,\Omega,(\mathcal{M},d),\Lambda \bigr)$ be an SCI computational problem with countable
$\Lambda=\{f_n\}_{n\in\NN}$. Let $I_\Lambda \subseteq \mathbb C^{\NN}$ be the
represented information space with the subspace product representation, and let
$\widehat{\Xi}:I_\Lambda \to \mathcal{M}$ be the factor map.

Fix also a representation $\delta_{\mathcal{M}}:\subseteq \mathbb N^{\mathbb N}\to \mathcal{M}$ and regard
$\widehat{\Xi}$ as a single-valued problem between represented spaces
\[
\widehat{\Xi}:\subseteq (I_\Lambda,\delta_{I_\Lambda}) \to (\mathcal{M},\delta_{\mathcal{M}}).
\]

Let $\mathcal R$ be a class of partial functionals
$\subseteq \mathbb N^{\mathbb N}\to \mathbb N^{\mathbb N}$.
For $k\in\mathbb N_0$ we say that $\mathcal{P}$ has a \textit{pure $\mathcal R$-$\Lim$-tower of
height $k$} if there exists $K\in\mathcal R$ such that
\[
\Limi{k} \circ K \preceq \mathcal N_{\widehat{\Xi}}.
\]
Equivalently: for every input name $p\in\dom(\delta_{I_\Lambda})$ naming some
$x\in I_\Lambda$, the name $K(p)$ lies in $\dom(\Limi{k})$ and
\[
\delta_{\mathcal{M}}\bigl( \Limi{k}(K(p)) \bigr)=\widehat{\Xi}(x).
\]

Define
\[
\mathrm{SCI}^{\Lim}_\mathcal R(\mathcal P) :=
\min\{k\in\mathbb N_0:
\mathcal P\text{ has a pure $\mathcal R-\Lim$ tower of height }k\},
\]
with value \(\infty\) if the set is empty.
\end{definition}

\begin{theorem}[Bridge theorem: pure $\mathcal R$-towers coincide with $\mathcal R$-Weihrauch rank]\label{thm:equivalence}
Assume the setting of \cref{def:SCIAR-package}. Assume moreover that the chosen
class $\mathcal R$ has the following \textit{$\Lim$-normal-form property}:
for every single-valued problem $f$ between represented spaces and every
$k\in\mathbb N_0$,
\[
f \le_{\mathrm W}^{\mathcal R} \Limi{k}
\quad\Longleftrightarrow\quad
\exists K\in\mathcal R \text{ such that } \Limi{k} \circ K \preceq \mathcal N_f.
\tag{$\ast_{\mathcal R,k}$}
\]
Then
\[
\mathrm{SCI}^{\Lim}_\mathcal R (\mathcal P) = \mathrm{rank}_{\mathcal R}(\widehat{\Xi}).
\]

In particular, for every class $\mathcal R$ for which $(\ast_{\mathcal R,k})$ holds for
all $k$, the pure $\mathcal R$-$\Lim$-tower height is exactly the $\mathcal R$-Weihrauch rank along the
iterated-limit chain.
\end{theorem}
\begin{proof}
By \cref{def:R-lim-rank},
\[
\mathrm{rank}_{\mathcal R}(\widehat{\Xi})
=
\min\Bigl\{ k\in\mathbb N_0 :
\widehat{\Xi} \le_{\mathrm W}^{\mathcal R} \Limi{k} \Bigr\},
\]
with value $\infty$ if the set is empty.

By the assumed $\Lim$-normal-form property $(\ast_{\mathcal R,k})$ (applied to $f=\widehat{\Xi}$), for each
$k\in\mathbb N_0$ we have
\[
\widehat{\Xi} \le_{\mathrm W}^{\mathcal R} \Limi{k}
\quad\Longleftrightarrow\quad
\exists K\in\mathcal R\text{ such that } \Limi{k} \circ K \preceq \mathcal N_{\widehat{\Xi}}.
\]
By \cref{def:SCIAR-package}, the right-hand side is exactly the statement that
$\mathcal{P}$ has a pure $\mathcal R$-$\Lim$-tower of height $k$.
Hence the set of admissible heights in the definition of
$\mathrm{SCI}^{\Lim}_\mathcal R(\mathcal{P})$ coincides with the set of admissible heights in the
definition of $\mathrm{rank}_{\mathcal R}(\widehat{\Xi})$.
Therefore the two minima coincide (or are both $\infty$).
\end{proof}

\subsection{Minimality Calculus}\label{sec:minimality-calculus}
\begin{theorem}[A $\mathsf{ZFC}$-minimality calculus]\label{thm:minimality}
Fix a universe of SCI problems $\mathcal P=\bigl( \Xi,\Omega,(\mathcal M,d),\Lambda \bigr)$ with countable
evaluation sets $\Lambda$.

Assume one aims for an \textit{effective} reading of the $\mathrm{SCI}_A$ interface satisfying the following
desiderata
\begin{enumerate}
\item[(D1)] \textbf{Computability-model stability:}
the admissible class of base post-processing operations forms a computability model
in the sense of \cref{def:compmodel} (well-defined morphism sets, identities, closure under
composition, and basic data operations).

\item[(D2)] \textbf{Weihrauch comparability via a pure $\mathcal R$-tower semantics:}
there exists a class $\mathcal R$ and an associated reducibility
$\le_{\mathrm W}^{\mathcal R}$ such that for every countable-$\Lambda$
SCI computational problem $\mathcal{P}$ with induced represented-space factor map
$\widehat{\Xi}:I_\Lambda \to \mathcal{M}$ one has
\[
\mathrm{SCI}^{\Lim}_\mathcal R(\mathcal{P})=\mathrm{rank}_{\mathcal R}(\widehat{\Xi}),
\]
where $\mathrm{SCI}^{\Lim}_\mathcal R(\mathcal{P})$ is the pure tower height from
\cref{def:SCIAR-package} and $\mathrm{rank}_{\mathcal R}$ is the
$\mathcal R$-Weihrauch rank along the iterated limit from \cref{def:R-lim-rank}.

\item[(D3)] \textbf{Nontriviality:}
there exists at least one problem $\mathcal P$ in the universe such that
$\mathrm{SCI}^{\Lim}_\mathcal R(\mathcal P)>0$.
\end{enumerate}

Consider then the following premises
\begin{enumerate}
\item[(P1)] \textbf{Fixed representations (up to a fixed equivalence class):}
the relevant input/output types are treated as represented spaces, i.e.\ the framework fixes
representations (or at least fixes admissible equivalence classes of representations) for the
data types involved in the computation model.

\item[(P2)] \textbf{Names, not atoms (when $\mathcal R$ is computability-based):}
if $\mathcal R$ is intended to capture computability on $\NN^\NN$ (e.g.\ $\mathsf{Comp}$ or $\mathsf{FMC}$),
then real/complex evaluation values must be supplied to post-processing by \textit{names}
(e.g.\ standard Cauchy names), not as atomic reals/complex numbers equipped with exact comparison
primitives.

\item[(P3)] \textbf{$\mathcal R$-restricted base post-processing:}
admissible base post-processing must be restricted to a distinguished class $\mathcal R$
(e.g.\ by requiring base maps to be $\mathcal R$-realizable on names).

\item[(P4)] \textbf{Adequacy of $\mathcal R$:}
the class $\mathcal R$ must be adequate (contain pairing/projections and constants and be closed
under composition) so that the resulting $\mathcal R$-reducibilities are at least preorders.

\item[(P5)] \textbf{$\Lim$-normal-form premise:}
for the chosen class $\mathcal R$, every reduction to $\Limi{k}$
can be put into pure pre-processing normal form, i.e.\
for every single-valued represented-space problem $f$ and every
$k\in\mathbb N_0$,
\[
f \le_{\mathrm W}^{\mathcal R} \Limi{k}
\quad\Longleftrightarrow\quad
\exists K\in\mathcal R \text{ such that } \Limi{k} \circ K \preceq \mathcal N_f.
\]
\end{enumerate}
Here, (P1)-(P4) are logically necessary for (D1)-(D3) and (P1)-(P5) are sufficient for (D1)-(D2).
\end{theorem}

\begin{proof}
(P1) is necessary by \cref{cor:need-representations}.
(P2) is necessary by \cref{cor:names-not-atoms}.
(P3) is necessary by \cref{thm:unrestricted-collapses}.
(P4) is necessary by \cref{lem:R-needed}.
(P1)-(P4) give a well-defined represented-space computability setting with a
$\mathcal R$-reduction preorder; then (D2) is exactly \cref{thm:equivalence}. 
More explicitly, under $(P1)$ the relevant data types are fixed represented spaces, so the
class of admissible $\mathcal R$-realizable maps is well-defined. By $(P4)$ this class contains
identities, pairing and projection maps, and constants, and it is closed under composition.
Hence the admissible base post-processing maps form a computability model in the sense
of \cref{def:compmodel}, i.e.\ $(D1)$ holds. Under $(P5)$, \cref{thm:equivalence} yields $(D2)$.
\end{proof}

\begin{remark}[Meta-level meaning of \enquote{logically necessary} and background logic]\label{rem:logical-necessity}
Throughout this work we reason in ordinary classical mathematics (e.g.\ a standard set-theoretic
metatheory such as $\mathsf{ZFC}$, with classical first-order logic in the background).
In particular, when we say that a premise $(P)$ is \textit{logically necessary} for a list of
desiderata $(D)$, we mean the following semantic implication:

\smallskip
\centerline{$(D)\ \Longrightarrow\ (P)$}

\smallskip
i.e.\ in the ambient metatheory there is \textit{no} interpretation of the $\mathrm{SCI}_A$ interface
satisfying $(D)$ while failing $(P)$ (equivalently: dropping $(P)$ admits a counterexample in which
at least one item of $(D)$ fails).
This is a \textit{semantic minimality} statement (in the sense of consequence/entailment),
not a reverse-mathematical claim about the minimal \textit{axiom system} in which the theorem can be proved.
If one wishes, all statements in \cref{sec:uniformity-minimality} can in principle be formalized in subsystems of analysis
(e.g.\ second-order arithmetic, see e.g., \cite{simpson09secAr}) after fixing codings of represented spaces, realizers, and
function classes; we do not pursue proof-theoretic strength here.
\end{remark}

\begin{remark}[High-level explanation tied to SCI literature]
\Cref{thm:minimality} formalizes:
type-$G$ (unrestricted) is not captured by Weihrauch degrees, hence one should restrict base post-processing
(continuous/Borel/Baire) and then compare to corresponding Weihrauch reducibilities.
For type-A, once one fixes representations/names, restricts base post-processing to
a class $\mathcal R$, and has the corresponding $\Lim$-normal-form property, the
\textit{pure} tower height agrees with the iterated-$\Lim$ rank.
\end{remark}

\section{Outlook}{\label{sec:outlook}}
Throughout this work we reason in ordinary classical mathematics, for instance in a standard set-theoretic metatheory such as $\mathsf{ZFC}$. The minimality statements in \cref{sec:uniformity-minimality} should be read as semantic minimality statements relative to that ambient metatheory: dropping one of the listed premises permits counterexamples to the desired computability-model or Weihrauch-comparability conclusions.

A reverse-mathematical analysis makes this dependence explicit by fixing a base theory
$T$ (typically $\mathsf{RCA}_0$) and internalizing the discussion as a \textit{scheme over codes}
for the relevant objects; see \cite{simpson09secAr} for a possible general framework.
Concretely, one restricts attention to instances of the SCI where the input/output types
are presented as represented spaces whose representations are coded by sets of naturals, and
where realizers/reductions are coded by indices or by accepted codes
for continuous/Borel functionals.  Under such a translation:
\begin{itemize}
\item (P1) becomes the requirement that representations (or admissible equivalence classes of
representations) are part of the data, so that \enquote{$\mathcal R$-computable/realizable} is a
\textit{definite} predicate on codes (cf.\ \cref{def:compmodel}(\textit{Well-definedness})).
\item (P2) becomes the requirement that the real/complex interface is name-based (e.g.\ Cauchy
names), not atomic exact comparison; this is precisely what makes computability imply
topological continuity on the induced representation-topologies (cf.\ the discontinuity obstruction in \cref{cor:names-not-atoms}).
\item (P3)-(P4) become closure axioms on the code-class for allowable pre-/post-processing,
ensuring that the induced reducibilities and height notions are stable under identity,
composition, and basic data operations (pairing/projections), i.e.\ that they define
preorders/hierarchies rather than ad hoc relations.
\end{itemize}
For type-A, after fixing representations/names and restricting admissible base
post-processing to a class $\mathcal R$, one obtains equality between pure tower
height and iterated-$\Lim$ rank precisely when the chosen $\mathcal R$ satisfies the
corresponding $\Lim$-normal-form property.

With this perspective, \cref{thm:minimality} suggests the following reverse-mathematics
questions, which are \textit{not} settled in the present work:
\begin{enumerate}
\item For computability-based choices of $\mathcal R$ (e.g.\ $\mathcal R=\Comp$ or
$\mathcal R=\mathsf{FMC}$), determine the weakest subsystem $T$ in which the full package
\[
(D1\wedge D2\wedge D3)\ \rightarrow\ (P1\wedge P2\wedge P3\wedge P4)
\]
can be formalized and proved.
\item For higher regularity choices of $\mathcal R$ (e.g.\ $\mathcal R =\mathsf{Cont}$ or $\mathcal R=\mathsf{Bor}$),
pin down how the answer depends on the chosen \textit{coding} of continuous/Borel functionals and on
the background axioms needed to develop the corresponding closure/measurability theory inside $T$
(e.g.\ potential needs for $\mathsf{WKL}_0$, $\mathsf{ACA}_0$, or $\mathsf{ATR}_0$-type principles).
\end{enumerate}
In short: \cref{thm:minimality} is a $\mathsf{ZFC}$-level minimality result, but it naturally opens the door to a
fine-grained calibration of its \textit{foundational strength} (and of the SCI-Weihrauch comparison theorems)
within the reverse-mathematics program.

\textbf{Acknowledgments } The author thanks Vasco Brattka, Matthias Schröder and Patrick Uftring for inspiring discussions around the different facets of this work.

\textbf{Statement } During the preparation of this work the author used UniBwM-ChatGPT5.2 in order to improve the language in the abstract and introduction. After using this tool, the author reviewed and edited the content as needed and takes full responsibility for the content of the published article.

\textbf{Conflict of interest } The author declares no conflict of interest.

% References
    \bibliographystyle{alpha}  
    \bibliography{bib.bib}
% ----------------------------------------

\end{document}

\typeout{get arXiv to do 4 passes: Label(s) may have changed. Rerun}
